\chapter{Relevés de la campagne de validation}
\label{ann:mesures}

Cette annexe regroupe l'intégralité des relevés de la campagne de validation, dans l'ordre où ils ont été effectués. Elle sert de source unique aux valeurs commentées dans le chapitre de résultats.

Chaque section suit la même organisation : un bloc de mesure par sous-section, avec son protocole immédiatement suivi de son tableau de relevés.

La colonne \textit{Valeur attendue} contient le critère fixé avant la mesure, à partir du schéma, des datasheets ou du dimensionnement. La colonne \textit{État} vaut OK lorsque la valeur mesurée respecte ce critère et \textbf{NOK} dans le cas contraire. Un point NOK n'est pas systématiquement un défaut de la carte, il signale que le critère initial n'a pas été vérifié tel qu'il avait été formulé. Les sept points NOK de la campagne sont tous commentés dans la colonne \textit{Remarques}.

\section{Traçabilité et bilan}
\label{ann:mesures_tracabilite}

Toutes les mesures matérielles et logicielles ont été réalisées sur la carte finale, entre le 27 et le 31 juillet 2026, avec le firmware de version \texttt{1c6e3c6} compilé sous ESP-IDF 6.0.1. Les instruments et les périphériques de référence n'ont pas changé pendant la campagne.

\begin{table}[H]
    \centering
    \caption{Moyens de mesure et périphériques de référence}
    \label{tab:ann_tracabilite}
    \small
    \begin{tabular}{@{}>{\raggedright\arraybackslash}p{0.38\textwidth} >{\raggedright\arraybackslash}p{0.42\textwidth}@{}}
        \toprule
        \textbf{Élément}                & \textbf{Référence}                       \\
        \midrule
        Multimètre et ohmmètre          & Keysight 34460A                          \\
        Alimentation de laboratoire     & GwInstek GPD-3303S                       \\
        Oscilloscope                    & Tektronix DPO 2014B                      \\
        Analyseur audio                 & Audio Precision ATS-2, logiciel ATS 1.60 \\
        Carte microSD                   & Swissbit 16 GB                           \\
        Casque, sortie jack             & Blon BL-03 32 $\Omega$                              \\
        Appareil Bluetooth de référence & JBL Charge 5                             \\
        \bottomrule
    \end{tabular}
\end{table}

La campagne compte 213 points de contrôle, dont 206 conformes au critère annoncé et 7 non conformes.

    {\small
        \setlength{\tabcolsep}{4pt}
        \renewcommand{\arraystretch}{1.1}
        \begin{longtable}{@{}>{\raggedright\arraybackslash}p{0.52\textwidth} r r r@{}}
            \caption{Bilan de la campagne, par bloc de mesure}
            \label{tab:ann_bilan}                                                                                 \\
            \toprule
            \textbf{Bloc}                                   & \textbf{Points} & \textbf{OK}  & \textbf{NOK}       \\
            \midrule
            \endfirsthead
            \caption[]{Bilan de la campagne, par bloc de mesure (suite)}                                          \\
            \toprule
            \textbf{Bloc}                                   & \textbf{Points} & \textbf{OK}  & \textbf{NOK}       \\
            \midrule
            \endhead
            \midrule
            \multicolumn{4}{r@{}}{\textit{Suite page suivante}}                                                   \\
            \endfoot
            \bottomrule
            \endlastfoot
            \multicolumn{4}{@{}l}{\textbf{Mesures matérielles, section~\ref{ann:mesures_hw}}}                     \\
            \addlinespace[0.2em]
            1. Vérification du \gls{pcb}, hors tension      & 11              & 11           & 0                  \\
            2. Alimentations sur batterie                   & 7               & 7            & 0                  \\
            3. Superviseur de tension XC6120                & 6               & 6            & 0                  \\
            4. Alimentations sur USB                        & 4               & 4            & 0                  \\
            5. \gls{strapping} et mux USB pour le flash     & 7               & 7            & 0                  \\
            6. Potentiomètre de volume                      & 3               & 3            & 0                  \\
            7. Bouton de reset matériel S1                  & 2               & 2            & 0                  \\
            8. Pull-ups et pull-downs de signaux            & 10              & 10           & 0                  \\
            9. Boutons                                      & 12              & 11           & 1                  \\
            10. Chaîne audio, alimentations et polarisation & 12              & 12           & 0                  \\
            11. Découpage du buck TPS62A01                  & 2               & 2            & 0                  \\
            12. Filtre RC différentiel                      & 4               & 4            & 0                  \\
            \midrule
            \multicolumn{4}{@{}l}{\textbf{Vérifications logicielles, section~\ref{ann:mesures_sw}}}               \\
            \addlinespace[0.2em]
            1. Menu principal                               & 8               & 8            & 0                  \\
            2. Statistics / Inputs                          & 3               & 3            & 0                  \\
            3. Statistics / Peripherals                     & 3               & 3            & 0                  \\
            4. Statistics / Storage                         & 5               & 5            & 0                  \\
            5. Statistics / Battery                         & 6               & 6            & 0                  \\
            6. Statistics / Battery / Debug                 & 2               & 2            & 0                  \\
            7. Storage                                      & 6               & 6            & 0                  \\
            8. Music                                        & 8               & 8            & 0                  \\
            9. Settings / Bluetooth                         & 6               & 6            & 0                  \\
            10. Settings / Audio                            & 2               & 2            & 0                  \\
            11. Settings / Screen                           & 2               & 2            & 0                  \\
            12. Settings / Power                            & 2               & 2            & 0                  \\
            13. Cas dégradés                                & 8               & 6            & 2                  \\
            14. Robustesse et durée                         & 3               & 3            & 0                  \\
            15. Absence de bruits parasites                 & 5               & 4            & 1                  \\
            \midrule
            \multicolumn{4}{@{}l}{\textbf{Formats audio, section~\ref{ann:mesures_formats}}}                      \\
            \addlinespace[0.2em]
            Formats supportés, sortie jack                  & 16              & 16           & 0                  \\
            Formats supportés, sortie Bluetooth             & 5               & 5            & 0                  \\
            \midrule
            \multicolumn{4}{@{}l}{\textbf{Caractérisation de la sortie filaire, section~\ref{ann:mesures_audio}}} \\
            \addlinespace[0.2em]
            A. Vérification du gain de la chaîne analogique & 6               & 6            & 0                  \\
            B. Nature de l'écrêtage à pleine échelle        & 11              & 8            & 3                  \\
            C. Préparation et câblage de la session         & 6               & 6            & 0                  \\
            D. Analyseur en configuration THD+N             & 11              & 11           & 0                  \\
            E. Analyseur en configuration \gls{diaphonie}   & 5               & 5            & 0                  \\
            F. Analyseur en configuration bruit             & 4               & 4            & 0                  \\
            \midrule
            \textbf{Total}                                  & \textbf{213}    & \textbf{206} & \textbf{7}         \\
        \end{longtable}
    }

\section{Mesures matérielles}
\label{ann:mesures_hw}

Le repérage des signaux suit les noms de nets du schéma. Aux blocs 1 à 4, le critère d'absence de court-circuit franc entre deux rails distincts est fixé à au moins 1 k$\Omega$, valeur qui sépare sans ambiguïté un défaut d'assemblage de la résistance d'isolement normale du montage.

\subsection{Bloc 1. Vérification du \gls{pcb}, hors tension}

Protocole : circuit hors tension, batterie déconnectée, mesures à l'ohmmètre.

{\scriptsize
    \setlength{\tabcolsep}{3pt}
    \renewcommand{\arraystretch}{1.15}
    \begin{longtable}{@{}
        >{\raggedright\arraybackslash}p{0.045\textwidth}
        >{\raggedright\arraybackslash}p{0.190\textwidth}
        >{\raggedright\arraybackslash}p{0.175\textwidth}
        >{\raggedright\arraybackslash}p{0.080\textwidth}
        >{\centering\arraybackslash}p{0.052\textwidth}
        >{\raggedright\arraybackslash}p{0.345\textwidth}@{}}
        \caption{Bloc 1 -- Vérification du \gls{pcb}, hors tension}
        \label{tab:ann_mesures_hw_b1} \\
        \toprule
        \textbf{N°} & \textbf{Point de mesure} & \textbf{Valeur attendue} & \textbf{Mesure} & \textbf{État} & \textbf{Remarques} \\
        \midrule
        \endfirsthead
        \caption[]{Bloc 1 -- Vérification du \gls{pcb}, hors tension (suite)} \\
        \toprule
        \textbf{N°} & \textbf{Point de mesure} & \textbf{Valeur attendue} & \textbf{Mesure} & \textbf{État} & \textbf{Remarques} \\
        \midrule
        \endhead
        \midrule
        \multicolumn{6}{r@{}}{\textit{Suite page suivante}} \\
        \endfoot
        \bottomrule
        \endlastfoot
        1.1  & \texttt{GND} - \texttt{GND}     & $<$ 0,2 $\Omega$  & 0,000 $\Omega$  & \OK & \\
        1.2  & \texttt{VBUS} - \texttt{GND}    & $\ge$ 1 k$\Omega$ & 44,1 k$\Omega$  & \OK & \\
        1.3  & \texttt{VSys} - \texttt{GND}    & $\ge$ 1 k$\Omega$ & 280 k$\Omega$   & \OK & \\
        1.4  & \texttt{VBatP} - \texttt{GND}   & $\ge$ 1 k$\Omega$ & 280 k$\Omega$   & \OK & \\
        1.5  & \texttt{+3.3V} - \texttt{GND}   & $\ge$ 1 k$\Omega$ & 1,15 k$\Omega$  & \OK & Valeur la plus basse de la série, cohérente avec les charges permanentes du rail \\
        1.6  & \texttt{VBUS} - \texttt{VSys}   & $\ge$ 1 k$\Omega$ & 21,1 M$\Omega$  & \OK & \\
        1.7  & \texttt{VBUS} - \texttt{+3.3V}  & $\ge$ 1 k$\Omega$ & 28 k$\Omega$    & \OK & \\
        1.8  & \texttt{VBUS} - \texttt{VBatP}  & $\ge$ 1 k$\Omega$ & 21,0 M$\Omega$  & \OK & \\
        1.9  & \texttt{VSys} - \texttt{+3.3V}  & $\ge$ 1 k$\Omega$ & 23,3 M$\Omega$  & \OK & \\
        1.10 & \texttt{VBatP} - \texttt{+3.3V} & $\ge$ 1 k$\Omega$ & 23,3 M$\Omega$  & \OK & \\
        1.11 & \texttt{VSys} - \texttt{VBatP}  & $<$ 0,2 $\Omega$  & 0,115 $\Omega$  & \OK & Les deux nets sont reliés par le chemin de puissance du chargeur \\
        \end{longtable}
        }

\subsection{Bloc 2. Alimentations sur batterie}

Protocole : alimentation de laboratoire en simulation de batterie, 3,5 V et 0,1 A, USB déconnecté.

{\scriptsize
\setlength{\tabcolsep}{3pt}
\renewcommand{\arraystretch}{1.15}
\begin{longtable}{@{}
>{\raggedright\arraybackslash}p{0.045\textwidth}
>{\raggedright\arraybackslash}p{0.190\textwidth}
>{\raggedright\arraybackslash}p{0.175\textwidth}
>{\raggedright\arraybackslash}p{0.080\textwidth}
>{\centering\arraybackslash}p{0.052\textwidth}
>{\raggedright\arraybackslash}p{0.345\textwidth}@{}}
\caption{Bloc 2 -- Alimentations sur batterie}
\label{tab:ann_mesures_hw_b2} \\
\toprule
\textbf{N°} & \textbf{Point de mesure} & \textbf{Valeur attendue} & \textbf{Mesure} & \textbf{État} & \textbf{Remarques} \\
        \midrule
        \endfirsthead
        \caption[]{Bloc 2 -- Alimentations sur batterie (suite)} \\
        \toprule
        \textbf{N°} & \textbf{Point de mesure} & \textbf{Valeur attendue} & \textbf{Mesure} & \textbf{État} & \textbf{Remarques} \\
        \midrule
        \endhead
        \midrule
        \multicolumn{6}{r@{}}{\textit{Suite page suivante}} \\
        \endfoot
        \bottomrule
        \endlastfoot
        2.1 & Courant d'alimentation          & $<$ 10 mA                                             & 1 mA    & \OK & \\
        2.2 & \texttt{VBatP} - \texttt{GND}   & 3,5 V                                                  & 3,495 V & \OK & \\
        2.3 & \texttt{VSys} - \texttt{GND}    & 3,5 V                                                  & 3,495 V & \OK & \\
        2.4 & \texttt{VBUS} - \texttt{GND}    & 0 V                                                    & 0,006 V & \OK & \\
        2.5 & \texttt{+3.3V} - \texttt{GND}   & 0 V                                                    & 0 V     & \OK & Sans appuyer sur le bouton A \\
        2.6 & \texttt{+3.3V} - \texttt{GND}   & 3,345 V, consigne fixée par R6 et R8                  & 3,345 V & \OK & En appuyant sur le bouton A \\
        2.7 & \texttt{BuckEN} - \texttt{GND}  & Niveau haut, $>$ 1,2 V                                & 3,345 V & \OK & En appuyant sur le bouton A \\
        \end{longtable}
        }

\subsection{Bloc 3. Superviseur de tension XC6120}

Protocole : alimentation de laboratoire abaissée par paliers, USB déconnecté.

{\scriptsize
\setlength{\tabcolsep}{3pt}
\renewcommand{\arraystretch}{1.15}
\begin{longtable}{@{}
>{\raggedright\arraybackslash}p{0.045\textwidth}
>{\raggedright\arraybackslash}p{0.190\textwidth}
>{\raggedright\arraybackslash}p{0.175\textwidth}
>{\raggedright\arraybackslash}p{0.080\textwidth}
>{\centering\arraybackslash}p{0.052\textwidth}
>{\raggedright\arraybackslash}p{0.345\textwidth}@{}}
\caption{Bloc 3 -- Superviseur de tension XC6120}
\label{tab:ann_mesures_hw_b3} \\
\toprule
\textbf{N°} & \textbf{Point de mesure} & \textbf{Valeur attendue} & \textbf{Mesure} & \textbf{État} & \textbf{Remarques} \\
        \midrule
        \endfirsthead
        \caption[]{Bloc 3 -- Superviseur de tension XC6120 (suite)} \\
        \toprule
        \textbf{N°} & \textbf{Point de mesure} & \textbf{Valeur attendue} & \textbf{Mesure} & \textbf{État} & \textbf{Remarques} \\
        \midrule
        \endhead
        \midrule
        \multicolumn{6}{r@{}}{\textit{Suite page suivante}} \\
        \endfoot
        \bottomrule
        \endlastfoot
        3.1 & \texttt{Supervisor} - \texttt{GND} & 3,25 V                                              & 3,244 V & \OK & Alimentation descendue à 3,25 V \\
        3.2 & \texttt{BuckEN} - \texttt{GND}     & Niveau haut, $>$ 1,2 V                              & 3,123 V & \OK & En appuyant sur le bouton A. Le superviseur n'a pas encore basculé \\
        3.3 & \texttt{Supervisor} - \texttt{GND} & 3,2 V                                                & 3,189 V & \OK & Alimentation descendue à 3,2 V \\
        3.4 & \texttt{BuckEN} - \texttt{GND}     & Niveau bas, $<$ 0,4 V                               & 0,017 V & \OK & En appuyant sur le bouton A. La protection contre la décharge profonde est fonctionnelle \\
        3.5 & Seuil de basculement                & Entre 3,136 V et 3,264 V \cite{torex_XC6120}        & 3,212 V & \OK & Balayage de l'alimentation de 3,40 V vers 3,10 V par pas de 10 mV, relevé de la tension exacte de bascule de \texttt{BuckEN} \\
        3.6 & Hystérésis                          & Tension de retour supérieure au seuil de bascule    & 3,279 V & \OK & Alimentation remontée, relevé de la tension de redémarrage \\
        \end{longtable}
        }

\subsection{Bloc 4. Alimentations sur USB}

Protocole : câble USB branché, alimentation de laboratoire débranchée, batterie déconnectée.

{\scriptsize
\setlength{\tabcolsep}{3pt}
\renewcommand{\arraystretch}{1.15}
\begin{longtable}{@{}
>{\raggedright\arraybackslash}p{0.045\textwidth}
>{\raggedright\arraybackslash}p{0.190\textwidth}
>{\raggedright\arraybackslash}p{0.175\textwidth}
>{\raggedright\arraybackslash}p{0.080\textwidth}
>{\centering\arraybackslash}p{0.052\textwidth}
>{\raggedright\arraybackslash}p{0.345\textwidth}@{}}
\caption{Bloc 4 -- Alimentations sur USB}
\label{tab:ann_mesures_hw_b4} \\
\toprule
\textbf{N°} & \textbf{Point de mesure} & \textbf{Valeur attendue} & \textbf{Mesure} & \textbf{État} & \textbf{Remarques} \\
        \midrule
        \endfirsthead
        \caption[]{Bloc 4 -- Alimentations sur USB (suite)} \\
        \toprule
        \textbf{N°} & \textbf{Point de mesure} & \textbf{Valeur attendue} & \textbf{Mesure} & \textbf{État} & \textbf{Remarques} \\
        \midrule
        \endhead
        \midrule
        \multicolumn{6}{r@{}}{\textit{Suite page suivante}} \\
        \endfoot
        \bottomrule
        \endlastfoot
        4.1 & \texttt{VBUS} - \texttt{GND}  & 5,0 à 5,5 V & 5,256 V & \OK & \\
        4.2 & \texttt{VBatP} - \texttt{GND} & 4,2 V       & 4,194 V & \OK & Tension imposée par le MAX77757 en l'absence de batterie \cite{analog_MAX77757} \\
        4.3 & \texttt{VSys} - \texttt{GND}  & 4,2 V       & 4,194 V & \OK & \\
        4.4 & \texttt{+3.3V} - \texttt{GND} & 3,345 V     & 3,339 V & \OK & \\
        \end{longtable}
        }

\subsection{Bloc 5. \gls{strapping} et mux USB pour le flash}

Protocole : USB branché, rail \texttt{+3.3V} actif. Seuils du module ESP32-WROVER-E, VIL = 0,825 V et VIH = 2,475 V \cite{ESP32_WROVER-E_Datasheet}.

{\scriptsize
\setlength{\tabcolsep}{3pt}
\renewcommand{\arraystretch}{1.15}
\begin{longtable}{@{}
>{\raggedright\arraybackslash}p{0.045\textwidth}
>{\raggedright\arraybackslash}p{0.190\textwidth}
>{\raggedright\arraybackslash}p{0.175\textwidth}
>{\raggedright\arraybackslash}p{0.080\textwidth}
>{\centering\arraybackslash}p{0.052\textwidth}
>{\raggedright\arraybackslash}p{0.345\textwidth}@{}}
\caption{Bloc 5 -- \gls{strapping} et mux USB pour le flash}
\label{tab:ann_mesures_hw_b5} \\
\toprule
\textbf{N°} & \textbf{Point de mesure} & \textbf{Valeur attendue} & \textbf{Mesure} & \textbf{État} & \textbf{Remarques} \\
        \midrule
        \endfirsthead
        \caption[]{Bloc 5 -- \gls{strapping} et mux USB pour le flash (suite)} \\
        \toprule
        \textbf{N°} & \textbf{Point de mesure} & \textbf{Valeur attendue} & \textbf{Mesure} & \textbf{État} & \textbf{Remarques} \\
        \midrule
        \endhead
        \midrule
        \multicolumn{6}{r@{}}{\textit{Suite page suivante}} \\
        \endfoot
        \bottomrule
        \endlastfoot
        5.1 & \texttt{fIO0}                  & Niveau haut, $>$ 2,475 V & 3,335 V & \OK & IO0, sélectionne le boot SPI \\
        5.2 & \texttt{DispRST}               & Niveau bas, $<$ 0,825 V  & 0,5 mV  & \OK & IO2, participe au mode de boot \\
        5.3 & \texttt{SPI.DC}                & Niveau bas, $<$ 0,825 V  & 0,5 mV  & \OK & IO12 (MTDI), fixe VDD\_SDIO à 3,3 V. Un niveau haut au reset basculerait la flash et la \gls{psram} sur le LDO 1,8 V interne et rendrait le module inopérant \cite{ESP32_Datasheet} \\
        5.4 & \texttt{SPI.CS}                & Niveau haut, $>$ 2,475 V & 3,329 V & \OK & IO15 (MTDO), log de boot actif \\
        5.5 & \texttt{SD.CS}                 & Niveau haut, $>$ 2,475 V & 3,336 V & \OK & IO5, timing SDIO slave \\
        5.6 & \texttt{fEN}                   & Niveau haut, $>$ 2,475 V & 3,335 V & \OK & EN, module hors reset \\
        5.7 & \texttt{USBDIR} - \texttt{GND} & Niveau bas, environ 0 V  & 0,001 V & \OK & Mux orienté vers le CP2102N, condition nécessaire au flash \\
        \end{longtable}
        }

\subsection{Bloc 6. Potentiomètre de volume}

Protocole : USB branché, rail \texttt{+3.3V} actif.

{\scriptsize
\setlength{\tabcolsep}{3pt}
\renewcommand{\arraystretch}{1.15}
\begin{longtable}{@{}
>{\raggedright\arraybackslash}p{0.045\textwidth}
>{\raggedright\arraybackslash}p{0.190\textwidth}
>{\raggedright\arraybackslash}p{0.175\textwidth}
>{\raggedright\arraybackslash}p{0.080\textwidth}
>{\centering\arraybackslash}p{0.052\textwidth}
>{\raggedright\arraybackslash}p{0.345\textwidth}@{}}
\caption{Bloc 6 -- Potentiomètre de volume}
\label{tab:ann_mesures_hw_b6} \\
\toprule
\textbf{N°} & \textbf{Point de mesure} & \textbf{Valeur attendue} & \textbf{Mesure} & \textbf{État} & \textbf{Remarques} \\
        \midrule
        \endfirsthead
        \caption[]{Bloc 6 -- Potentiomètre de volume (suite)} \\
        \toprule
        \textbf{N°} & \textbf{Point de mesure} & \textbf{Valeur attendue} & \textbf{Mesure} & \textbf{État} & \textbf{Remarques} \\
        \midrule
        \endhead
        \midrule
        \multicolumn{6}{r@{}}{\textit{Suite page suivante}} \\
        \endfoot
        \bottomrule
        \endlastfoot
        6.1 & \texttt{Vol} - \texttt{GND} & 0,147 V $\pm$ 30 \% & 0,141 V & \OK & Potentiomètre R13 en butée basse. Bornes calculées pour rester dans la plage linéaire de l'\gls{adc} \\
        6.2 & \texttt{Vol} - \texttt{GND} & 2,518 V $\pm$ 30 \% & 2,582 V & \OK & Potentiomètre R13 en butée haute \\
        6.3 & \texttt{Vol} - \texttt{GND} & Valeur intermédiaire entre 6.1 et 6.2 & 1,380 V & \OK & Position médiane. Vérifie l'absence de discontinuité de la piste \\
        \end{longtable}
        }

\subsection{Bloc 7. Bouton de reset matériel S1}

Protocole : USB branché, appui maintenu sur S1.

{\scriptsize
\setlength{\tabcolsep}{3pt}
\renewcommand{\arraystretch}{1.15}
\begin{longtable}{@{}
>{\raggedright\arraybackslash}p{0.045\textwidth}
>{\raggedright\arraybackslash}p{0.190\textwidth}
>{\raggedright\arraybackslash}p{0.175\textwidth}
>{\raggedright\arraybackslash}p{0.080\textwidth}
>{\centering\arraybackslash}p{0.052\textwidth}
>{\raggedright\arraybackslash}p{0.345\textwidth}@{}}
\caption{Bloc 7 -- Bouton de reset matériel S1}
\label{tab:ann_mesures_hw_b7} \\
\toprule
\textbf{N°} & \textbf{Point de mesure} & \textbf{Valeur attendue} & \textbf{Mesure} & \textbf{État} & \textbf{Remarques} \\
        \midrule
        \endfirsthead
        \caption[]{Bloc 7 -- Bouton de reset matériel S1 (suite)} \\
        \toprule
        \textbf{N°} & \textbf{Point de mesure} & \textbf{Valeur attendue} & \textbf{Mesure} & \textbf{État} & \textbf{Remarques} \\
        \midrule
        \endhead
        \midrule
        \multicolumn{6}{r@{}}{\textit{Suite page suivante}} \\
        \endfoot
        \bottomrule
        \endlastfoot
        7.1 & \texttt{BuckEN} - \texttt{GND} & Niveau bas, $<$ 0,4 V & 0,001 V & \OK & \\
        7.2 & \texttt{+3.3V} - \texttt{GND}  & 0 V                   & 0,003 V & \OK & Le rail s'effondre, le force-reset matériel est fonctionnel \\
        \end{longtable}
        }

\subsection{Bloc 8. Pull-ups et pull-downs de signaux}

Protocole : USB branché, rail \texttt{+3.3V} actif. Au \gls{por}, les I/O du PI4IOE5V9554 sont en entrée avec une pull-up interne de 100 k$\Omega$ typique vers VCC \cite{diodes_inc_PI4IOE5V9554}.

{\scriptsize
\setlength{\tabcolsep}{3pt}
\renewcommand{\arraystretch}{1.15}
\begin{longtable}{@{}
>{\raggedright\arraybackslash}p{0.045\textwidth}
>{\raggedright\arraybackslash}p{0.190\textwidth}
>{\raggedright\arraybackslash}p{0.175\textwidth}
>{\raggedright\arraybackslash}p{0.080\textwidth}
>{\centering\arraybackslash}p{0.052\textwidth}
>{\raggedright\arraybackslash}p{0.345\textwidth}@{}}
\caption{Bloc 8 -- Pull-ups et pull-downs de signaux}
\label{tab:ann_mesures_hw_b8} \\
\toprule
\textbf{N°} & \textbf{Point de mesure} & \textbf{Valeur attendue} & \textbf{Mesure} & \textbf{État} & \textbf{Remarques} \\
        \midrule
        \endfirsthead
        \caption[]{Bloc 8 -- Pull-ups et pull-downs de signaux (suite)} \\
        \toprule
        \textbf{N°} & \textbf{Point de mesure} & \textbf{Valeur attendue} & \textbf{Mesure} & \textbf{État} & \textbf{Remarques} \\
        \midrule
        \endhead
        \midrule
        \multicolumn{6}{r@{}}{\textit{Suite page suivante}} \\
        \endfoot
        \bottomrule
        \endlastfoot
        8.1  & \texttt{BatAlrt} - \texttt{GND}  & Niveau haut & 3,334 V & \OK & Sortie ALRT du MAX17260, open-drain active basse \cite{analog_MAX17260} \\
        8.2  & \texttt{SDA} - \texttt{GND}      & Niveau haut & 3,336 V & \OK & Bus au repos, aucun esclave ne tire la ligne \\
        8.3  & \texttt{SCL} - \texttt{GND}      & Niveau haut & 3,336 V & \OK & \\
        8.4  & \texttt{DACMute} - \texttt{GND}  & Niveau bas, $<$ 0,3 $\times$ DVDD = 1,00 V \cite{TI_PCM5242} & 0,353 V & \OK & Pont entre la pull-up interne de 100 k$\Omega$ de l'expanseur et la pull-down externe de 10 k$\Omega$, soit 3,345 $\times$ 10 / 110 = 0,304 V. Le \gls{DAC} est muet au \gls{por} \\
        8.5  & \texttt{AmpSHDN} - \texttt{GND}  & Niveau bas, $<$ 0,4 V & 0,361 V & \OK & Même pont diviseur qu'au point 8.4 \\
        8.6  & \texttt{SD.CS}                   & Niveau haut & 3,336 V & \OK & Doublon du point 5.5, conservé pour la cohérence du bloc \\
        8.7  & \texttt{SD.MOSI}                 & Niveau haut & 3,336 V & \OK & \\
        8.8  & \texttt{SD.MISO}                 & Niveau haut & 3,336 V & \OK & \\
        8.9  & \texttt{SD.Detect}               & Niveau haut & 3,336 V & \OK & Aucune carte insérée \\
        8.10 & \texttt{SD.Detect}               & Niveau bas  & 0,001 V & \OK & Carte microSD insérée \\
        \end{longtable}
        }

\subsection{Bloc 9. Boutons}

Protocole : USB branché, rail \texttt{+3.3V} actif. Seuils du PI4IOE5V9554, VIL = 0,81 V et VIH = 1,8 V \cite{diodes_inc_PI4IOE5V9554}. Seuils de l'ESP32 alimenté en 3,3 V, VIL = 0,825 V et VIH = 2,475 V \cite{ESP32_WROVER-E_Datasheet}.

{\scriptsize
\setlength{\tabcolsep}{3pt}
\renewcommand{\arraystretch}{1.15}
\begin{longtable}{@{}
>{\raggedright\arraybackslash}p{0.045\textwidth}
>{\raggedright\arraybackslash}p{0.190\textwidth}
>{\raggedright\arraybackslash}p{0.175\textwidth}
>{\raggedright\arraybackslash}p{0.080\textwidth}
>{\centering\arraybackslash}p{0.052\textwidth}
>{\raggedright\arraybackslash}p{0.345\textwidth}@{}}
\caption{Bloc 9 -- Boutons}
\label{tab:ann_mesures_hw_b9} \\
\toprule
\textbf{N°} & \textbf{Point de mesure} & \textbf{Valeur attendue} & \textbf{Mesure} & \textbf{État} & \textbf{Remarques} \\
        \midrule
        \endfirsthead
        \caption[]{Bloc 9 -- Boutons (suite)} \\
        \toprule
        \textbf{N°} & \textbf{Point de mesure} & \textbf{Valeur attendue} & \textbf{Mesure} & \textbf{État} & \textbf{Remarques} \\
        \midrule
        \endhead
        \midrule
        \multicolumn{6}{r@{}}{\textit{Suite page suivante}} \\
        \endfoot
        \bottomrule
        \endlastfoot
        9.1  & \texttt{bA} - \texttt{GND} & Niveau bas, $<$ 0,825 V  & 0,009 V & \OK  & Bouton au repos \\
        9.2  & \texttt{bA} - \texttt{GND} & Niveau haut, $>$ 2,475 V & 2,221 V & \NOK & Bouton appuyé. Le bouton A est référencé à \texttt{VSys} à travers une diode Zener de 3,3 V et non au rail \texttt{+3.3V}, car il doit agir rail éteint. Le point 2.1 de la section~\ref{ann:mesures_sw} montre que la détection reste fiable \\
        9.3  & \texttt{bB} - \texttt{GND} & Niveau haut, $>$ 1,8 V   & 3,336 V & \OK  & Bouton au repos \\
        9.4  & \texttt{bB} - \texttt{GND} & Niveau bas, $<$ 0,81 V   & 0,000 V & \OK  & Bouton appuyé \\
        9.5  & \texttt{bU} - \texttt{GND} & Niveau haut, $>$ 1,8 V   & 3,336 V & \OK  & Bouton au repos \\
        9.6  & \texttt{bU} - \texttt{GND} & Niveau bas, $<$ 0,81 V   & 0,000 V & \OK  & Bouton appuyé \\
        9.7  & \texttt{bD} - \texttt{GND} & Niveau haut, $>$ 1,8 V   & 3,336 V & \OK  & Bouton au repos \\
        9.8  & \texttt{bD} - \texttt{GND} & Niveau bas, $<$ 0,81 V   & 0,000 V & \OK  & Bouton appuyé \\
        9.9  & \texttt{bL} - \texttt{GND} & Niveau haut, $>$ 1,8 V   & 3,336 V & \OK  & Bouton au repos \\
        9.10 & \texttt{bL} - \texttt{GND} & Niveau bas, $<$ 0,81 V   & 0,000 V & \OK  & Bouton appuyé \\
        9.11 & \texttt{bR} - \texttt{GND} & Niveau haut, $>$ 1,8 V   & 3,336 V & \OK  & Bouton au repos \\
        9.12 & \texttt{bR} - \texttt{GND} & Niveau bas, $<$ 0,81 V   & 0,000 V & \OK  & Bouton appuyé \\
        \end{longtable}
        }

\subsection{Bloc 10. Chaîne audio, alimentations et polarisation}

Protocole : USB branché, rail \texttt{+3.3V} actif, casque débranché. Les points 10.9 et 10.10 sont relevés pendant la lecture d'un fichier à 1 kHz, les autres au repos. Repérage des pins selon la feuille P8 du schéma. Mesures au multimètre. L'amplificateur MAX97220A emploie une topologie \gls{directdrive}, qui supprime les condensateurs de liaison et rend l'offset directement mesurable au connecteur \cite{analog_MAX97220A}.

{\scriptsize
\setlength{\tabcolsep}{3pt}
\renewcommand{\arraystretch}{1.15}
\begin{longtable}{@{}
>{\raggedright\arraybackslash}p{0.045\textwidth}
>{\raggedright\arraybackslash}p{0.190\textwidth}
>{\raggedright\arraybackslash}p{0.175\textwidth}
>{\raggedright\arraybackslash}p{0.080\textwidth}
>{\centering\arraybackslash}p{0.052\textwidth}
>{\raggedright\arraybackslash}p{0.345\textwidth}@{}}
\caption{Bloc 10 -- Chaîne audio, alimentations et polarisation}
\label{tab:ann_mesures_hw_b10} \\
\toprule
\textbf{N°} & \textbf{Point de mesure} & \textbf{Valeur attendue} & \textbf{Mesure} & \textbf{État} & \textbf{Remarques} \\
        \midrule
        \endfirsthead
        \caption[]{Bloc 10 -- Chaîne audio, alimentations et polarisation (suite)} \\
        \toprule
        \textbf{N°} & \textbf{Point de mesure} & \textbf{Valeur attendue} & \textbf{Mesure} & \textbf{État} & \textbf{Remarques} \\
        \midrule
        \endhead
        \midrule
        \multicolumn{6}{r@{}}{\textit{Suite page suivante}} \\
        \endfoot
        \bottomrule
        \endlastfoot
        10.1  & AVDD du PCM5242, pin 14   & 3,345 V & 3,338 V & \OK & \\
        10.2  & DVDD du PCM5242, pin 4    & 3,345 V & 3,336 V & \OK & \\
        10.3  & CPVDD du PCM5242, pin 5   & 3,345 V & 3,336 V & \OK & Alimentation de la pompe de charge du \gls{DAC} \\
        10.4  & VCOM du PCM5242, pin 16   & 0 V     & 0,000 V & \OK & \\
        10.5  & PVDD du MAX97220A, pin 1  & 3,345 V & 3,336 V & \OK & Fixe le plafond d'écrêtage de l'étage de sortie. Valeur reprise au bloc B de la section~\ref{ann:mesures_audio} \\
        10.6  & SVDD du MAX97220A, pin 13 & 3,345 V & 3,336 V & \OK & \\
        10.7  & SVDD2 du MAX97220A, pin 9 & 3,345 V & 3,336 V & \OK & \\
        10.8  & PVSS du MAX97220A, pin 5  & Environ $-3,3$ V & $-3,303$ V & \OK & Rail négatif produit par la pompe de charge interne \cite{analog_MAX97220A} \\
        10.9  & XSMT du PCM5242, pin 1    & Niveau haut, $>$ 2,0 V, en lecture \cite{TI_PCM5242} & 3,333 V & \OK & Le \gls{DAC} n'est pas muté \\
        10.10 & SHDN du MAX97220A, pin 16 & Niveau haut, $>$ 1,4 V, en lecture \cite{analog_MAX97220A} & 3,333 V & \OK & L'amplificateur est actif \\
        10.11 & Jack tip - jack sleeve    & $<$ 350 $\mu$V & 55 $\mu$V & \OK & VOS maximal du MAX97220A \cite{analog_MAX97220A}. La topologie \gls{directdrive} supprime les condensateurs de liaison, l'offset est donc directement mesurable au connecteur \\
        10.12 & Jack ring - jack sleeve   & $<$ 350 $\mu$V & 55 $\mu$V & \OK & \\
        \end{longtable}
        }

\subsection{Bloc 11. Découpage du buck TPS62A01}

Protocole : oscilloscope Tektronix DPO 2014B, alimentation sur batterie, lecture d'un fichier à 1 kHz en cours. Le point 11.1 est relevé sur le nœud SW, pin 2. Le point 11.2 est relevé sur le rail \texttt{+3.3V}, sonde 10:1, pointe posée directement sur les pads de C10, couplage alternatif.

{\scriptsize
\setlength{\tabcolsep}{3pt}
\renewcommand{\arraystretch}{1.15}
\begin{longtable}{@{}
>{\raggedright\arraybackslash}p{0.045\textwidth}
>{\raggedright\arraybackslash}p{0.190\textwidth}
>{\raggedright\arraybackslash}p{0.175\textwidth}
>{\raggedright\arraybackslash}p{0.080\textwidth}
>{\centering\arraybackslash}p{0.052\textwidth}
>{\raggedright\arraybackslash}p{0.345\textwidth}@{}}
\caption{Bloc 11 -- Découpage du buck TPS62A01}
\label{tab:ann_mesures_hw_b11} \\
\toprule
\textbf{N°} & \textbf{Point de mesure} & \textbf{Valeur attendue} & \textbf{Mesure} & \textbf{État} & \textbf{Remarques} \\
        \midrule
        \endfirsthead
        \caption[]{Bloc 11 -- Découpage du buck TPS62A01 (suite)} \\
        \toprule
        \textbf{N°} & \textbf{Point de mesure} & \textbf{Valeur attendue} & \textbf{Mesure} & \textbf{État} & \textbf{Remarques} \\
        \midrule
        \endhead
        \midrule
        \multicolumn{6}{r@{}}{\textit{Suite page suivante}} \\
        \endfoot
        \bottomrule
        \endlastfoot
        11.1 & Fréquence de découpage sur SW, pin 2, en lecture             & 2,4 MHz typique en mode \gls{pwm} & 2,083 MHz & \OK & Le TPS62A0x n'est pas à fréquence fixe. Le temps de blocage est fixé par un circuit adaptatif à partir du rapport VIN sur VOUT, ce qui ne rend la fréquence qu'approximativement constante \cite{TI_TPS62A01} \\
        11.2 & Ondulation crête à crête du rail \texttt{+3.3V}, en lecture & Quelques millivolts                & 9,60 mV   & \OK & Soit 0,29 \% du rail \\
        \end{longtable}
        }

\subsection{Bloc 12. Filtre RC différentiel}

Protocole : oscilloscope Tektronix DPO 2014B, deux voies acquises simultanément avec des sondes 10:1 identiques et compensées, couplage alternatif, référence sur \texttt{GND}. CH1 en amont du filtre sur \texttt{OUTLP}, CH2 en aval sur \texttt{INR+}. Le croisement des voies L et R au niveau de l'amplificateur est celui documenté sur la feuille P8 du schéma, \texttt{OUTLP} attaquant bien l'entrée non inverseuse. Les points 12.1 à 12.4 exploitent le rapport CH2 sur CH1, mesuré par la fonction Cycle \gls{rms} sur les deux voies simultanément, calibres verticaux inchangés sur toute la série. Fichiers de test produits par \texttt{gen\_fichiers\_test\_ap.py}, jeu filtre, à 192 kHz et $-6$ \gls{dbfs}.

{\scriptsize
\setlength{\tabcolsep}{3pt}
\renewcommand{\arraystretch}{1.15}
\begin{longtable}{@{}
>{\raggedright\arraybackslash}p{0.045\textwidth}
>{\raggedright\arraybackslash}p{0.190\textwidth}
>{\raggedright\arraybackslash}p{0.175\textwidth}
>{\raggedright\arraybackslash}p{0.080\textwidth}
>{\centering\arraybackslash}p{0.052\textwidth}
>{\raggedright\arraybackslash}p{0.345\textwidth}@{}}
\caption{Bloc 12 -- Filtre RC différentiel}
\label{tab:ann_mesures_hw_b12} \\
\toprule
\textbf{N°} & \textbf{Point de mesure} & \textbf{Valeur attendue} & \textbf{Mesure} & \textbf{État} & \textbf{Remarques} \\
        \midrule
        \endfirsthead
        \caption[]{Bloc 12 -- Filtre RC différentiel (suite)} \\
        \toprule
        \textbf{N°} & \textbf{Point de mesure} & \textbf{Valeur attendue} & \textbf{Mesure} & \textbf{État} & \textbf{Remarques} \\
        \midrule
        \endhead
        \midrule
        \multicolumn{6}{r@{}}{\textit{Suite page suivante}} \\
        \endfoot
        \bottomrule
        \endlastfoot
        12.1 & Fondamentale 1 kHz en amont du filtre, sur \texttt{OUTLP}                              & Environ 1,9 Vrms & 1,88 Vrms & \OK & Une seule branche, soit la moitié de l'amplitude différentielle de 3,74 Vrms. Sert de référence d'entrée à tous les rapports du bloc \\
        12.2 & Fondamentale 1 kHz en aval du filtre, sur \texttt{INR+}                                & Un tiers du point 12.1, soit $-9,54$ dB, le rapport du pont diviseur d'entrée de l'amplificateur valant RF / (RIN + RF) = 1/3 pour un gain de 0,5 & 0,62 Vrms & \OK & Le nœud \texttt{INR+} n'est pas la sortie du filtre mais un point interne du réseau d'entrée du MAX97220A. L'atténuation observée est plate et purement résistive, sans rapport avec le filtre \\
        12.3 & Transparence du filtre dans la bande audio, rapport \texttt{INR+} sur \texttt{OUTLP} de 1 kHz à 20 kHz & Écart inférieur à 0,1 dB & 0,01 dB & \OK & Série complète au tableau~\ref{tab:ann_filtre_rc} \\
        12.4 & Capture avant et après le filtre, balayage à 90 kHz                                    & Paliers de reconstruction nettement visibles sur \texttt{OUTLP}, lissés sur \texttt{INR+} & Conforme & \OK & Résultat principal du bloc, figure~\ref{filtre_rc_capture.PNG}. En amont du filtre les paliers de conversion du \gls{DAC} sont visibles, en aval le signal est lisse \\
    \end{longtable}
}

\fig[H, width=1\textwidth]{Signal avant et après le filtre RC, balayage à 90 kHz}{filtre_rc_capture.PNG}

\section{Transparence du filtre RC différentiel}
\label{ann:mesures_filtre}

Ces relevés viennent à l'appui du point 12.4 du tableau~\ref{tab:ann_mesures_hw_b12}. Leur objet est de vérifier que le filtre n'intervient pas dans la bande utile. Ils ne déterminent pas la fréquence de coupure, qui se situe très au-delà de la plage accessible avec des fichiers à 192 kHz.

Les fichiers WAV sont produits par \texttt{gen\_fichiers\_test\_ap.py}, dossier filtre, à 192 kHz et $-6$ \gls{dbfs}. Ils suivent la convention \texttt{SWEEP\_<f>kHz\_192k\_-6dBFS.wav}, la fréquence de la première colonne du tableau désignant donc sans ambiguïté le fichier joué. Les sondes 10:1 sont identiques et compensées, en couplage alternatif, référencées sur \texttt{GND}. CH1 est placée sur \texttt{OUTLP} et CH2 sur \texttt{INR+}, la mesure automatique Cycle \gls{rms} étant relevée sur les deux voies simultanément. Les calibres verticaux restent inchangés sur toute la série. Le rapport est normalisé par le relevé à 1 kHz, ce qui élimine à la fois l'écart de gain entre les deux voies de l'oscilloscope et l'atténuation plate du réseau d'entrée de l'amplificateur.

\begin{table}[H]
    \centering
    \caption{Rapport \texttt{INR+} sur \texttt{OUTLP} en fonction de la fréquence}
    \label{tab:ann_filtre_rc}
    \small
    \setlength{\tabcolsep}{6pt}
    \begin{tabular}{@{}r r r r r@{}}
        \toprule
        \textbf{f [kHz]} & \textbf{V(OUTLP) [mV]} & \textbf{V(INR+) [mV]} & \textbf{Rapport} & \textbf{Écart / 1 kHz [dB]} \\
        \midrule
        1                & 1060                   & 350                   & 0,3302           & 0,000                       \\
        20               & 1070                   & 353                   & 0,3299           & $-0,007$                    \\
        40               & 1080                   & 352                   & 0,3259           & $-0,113$                    \\
        60               & 1070                   & 350                   & 0,3271           & $-0,082$                    \\
        80               & 1070                   & 347                   & 0,3243           & $-0,156$                    \\
        90               & 1060                   & 345                   & 0,3255           & $-0,125$                    \\
        \bottomrule
    \end{tabular}
\end{table}

Le rapport reste constant à 0,33 sur toute la plage explorée. L'écart par rapport au relevé à 1 kHz vaut 0,01 dB à 20 kHz et n'excède jamais 0,16 dB jusqu'à 90 kHz, valeur du même ordre que la résolution de lecture de l'oscilloscope. Le filtre est donc transparent dans la bande audio et le reste largement au-delà, ce qui est l'exigence de conception recherchée.

Quatre limites doivent accompagner ces chiffres.

\begin{itemize}
    \item La résolution de lecture est de 10 mV sur environ 1070 mV en amont, soit 0,08 dB. Les écarts relevés au-delà de la bande audio sont du même ordre et ne doivent pas être interprétés comme une réponse en fréquence.
    \item Le rapport n'est pas monotone entre 40 et 60 kHz, ce qui confirme que la dispersion est dominée par la quantification de lecture et non par le comportement du filtre.
    \item La capacité d'entrée de la sonde, de 10 à 12 pF, s'ajoute au condensateur du filtre au nœud \texttt{INR+}. Elle reste sans effet mesurable dans la bande audio.
    \item Une mesure de déphasage a été tentée puis abandonnée. Les paliers de reconstruction du \gls{DAC}, à 2,1 points par période à 90 kHz, rendent la détection des passages par zéro instable.
\end{itemize}

\section{Vérifications logicielles}
\label{ann:mesures_sw}

Ces vérifications portent sur le firmware final et suivent l'arborescence de l'interface utilisateur. Elles sont fonctionnelles et se lisent comme une recette d'acceptation, chaque ligne énonçant le comportement attendu et le résultat observé. Sauf mention contraire dans le protocole d'un bloc, la carte est alimentée sur batterie, démarrée par le bouton A, et la carte microSD de référence est insérée.

\subsection{Bloc 1. Menu principal}

Protocole : démarrage par le bouton A, batterie chargée, carte microSD insérée.

{\footnotesize
\setlength{\tabcolsep}{4pt}
\renewcommand{\arraystretch}{1.15}
\begin{longtable}{@{}
    >{\raggedright\arraybackslash}p{0.045\textwidth}
    >{\raggedright\arraybackslash}p{0.400\textwidth}
    >{\centering\arraybackslash}p{0.055\textwidth}
    >{\raggedright\arraybackslash}p{0.390\textwidth}@{}}
    \caption{Bloc 1 -- Menu principal}
    \label{tab:ann_mesures_sw_b1}                                                                                                                      \\
    \toprule
    \textbf{N°} & \textbf{Comportement attendu}                                                                   & \textbf{État} & \textbf{Remarques} \\
    \midrule
    \endfirsthead
    \caption[]{Bloc 1 -- Menu principal (suite)}                                                                                                       \\
    \toprule
    \textbf{N°} & \textbf{Comportement attendu}                                                                   & \textbf{État} & \textbf{Remarques} \\
    \midrule
    \endhead
    \midrule
    \multicolumn{4}{r@{}}{\textit{Suite page suivante}}                                                                                                \\
    \endfoot
    \bottomrule
    \endlastfoot
    1.1         & Écran allumé et menu principal affiché après le démarrage                                       & \OK           &                    \\
    1.2         & Les couleurs primaires s'affichent sans glitch ni inversion                                     & \OK           &                    \\
    1.3         & La barre de statut affiche l'état de l'antenne, cohérent avec l'état réel du Bluetooth          & \OK           &                    \\
    1.4         & La barre de statut affiche la sortie active et suit le réglage de Settings / Audio              & \OK           &                    \\
    1.5         & Le pourcentage affiché correspond au RepSOC lu dans Statistics / Battery                        & \OK           &                    \\
    1.6         & Les quatre directions pré-sélectionnent les quatre menus, la sélection reboucle en fin de liste & \OK           &                    \\
    1.7         & Le bouton A ouvre le menu pré-sélectionné                                                       & \OK           &                    \\
    1.8         & Le bouton B revient au menu précédent depuis chacun des quatre menus                            & \OK           &                    \\
\end{longtable}
}

\subsection{Bloc 2. Statistics / Inputs}

{\footnotesize
    \setlength{\tabcolsep}{4pt}
    \renewcommand{\arraystretch}{1.15}
    \begin{longtable}{@{}
        >{\raggedright\arraybackslash}p{0.045\textwidth}
        >{\raggedright\arraybackslash}p{0.400\textwidth}
        >{\centering\arraybackslash}p{0.055\textwidth}
        >{\raggedright\arraybackslash}p{0.390\textwidth}@{}}
        \caption{Bloc 2 -- Statistics / Inputs}
        \label{tab:ann_mesures_sw_b2} \\
        \toprule
        \textbf{N°} & \textbf{Comportement attendu} & \textbf{État} & \textbf{Remarques} \\
        \midrule
        \endfirsthead
        \caption[]{Bloc 2 -- Statistics / Inputs (suite)} \\
        \toprule
        \textbf{N°} & \textbf{Comportement attendu} & \textbf{État} & \textbf{Remarques} \\
        \midrule
        \endhead
        \midrule
        \multicolumn{4}{r@{}}{\textit{Suite page suivante}} \\
        \endfoot
        \bottomrule
        \endlastfoot
        2.1 & Chaque appui incrémente le compteur du bouton concerné de 1, et d'aucun autre & \OK & Malgré la faible tension mesurée sur \texttt{bA} au point 9.2 du tableau~\ref{tab:ann_mesures_hw_b9}, l'ESP32 détecte l'appui de manière répétitive \\
        2.2 & Potentiomètre immobile pendant 60 s, la valeur affichée ne dérive pas          & \OK & La valeur dérive de $\pm$ 3 mV. Le volume fonctionne par paliers et une dérive de cette amplitude ne suffit pas à franchir un palier \\
        2.3 & En butée basse puis en butée haute, les valeurs affichées encadrent la plage attendue & \OK & \\
        \end{longtable}
        }

\subsection{Bloc 3. Statistics / Peripherals}

{\footnotesize
\setlength{\tabcolsep}{4pt}
\renewcommand{\arraystretch}{1.15}
\begin{longtable}{@{}
>{\raggedright\arraybackslash}p{0.045\textwidth}
>{\raggedright\arraybackslash}p{0.400\textwidth}
>{\centering\arraybackslash}p{0.055\textwidth}
>{\raggedright\arraybackslash}p{0.390\textwidth}@{}}
\caption{Bloc 3 -- Statistics / Peripherals}
\label{tab:ann_mesures_sw_b3} \\
\toprule
\textbf{N°} & \textbf{Comportement attendu} & \textbf{État} & \textbf{Remarques} \\
        \midrule
        \endfirsthead
        \caption[]{Bloc 3 -- Statistics / Peripherals (suite)} \\
        \toprule
        \textbf{N°} & \textbf{Comportement attendu} & \textbf{État} & \textbf{Remarques} \\
        \midrule
        \endhead
        \midrule
        \multicolumn{4}{r@{}}{\textit{Suite page suivante}} \\
        \endfoot
        \bottomrule
        \endlastfoot
        3.1 & Les registres relus sur le PCM5242, le MAX17260 et l'expanseur correspondent aux valeurs écrites & \OK & Valide au passage les adresses \gls{i2c} 0x4C, 0x36 et 0x38 \\
        3.2 & L'état de montage de la carte SD suit l'insertion et le retrait                                  & \OK & \\
        3.3 & L'état du Bluetooth suit la connexion et la déconnexion                                          & \OK & \\
        \end{longtable}
        }

\subsection{Bloc 4. Statistics / Storage}

{\footnotesize
\setlength{\tabcolsep}{4pt}
\renewcommand{\arraystretch}{1.15}
\begin{longtable}{@{}
>{\raggedright\arraybackslash}p{0.045\textwidth}
>{\raggedright\arraybackslash}p{0.400\textwidth}
>{\centering\arraybackslash}p{0.055\textwidth}
>{\raggedright\arraybackslash}p{0.390\textwidth}@{}}
\caption{Bloc 4 -- Statistics / Storage}
\label{tab:ann_mesures_sw_b4} \\
\toprule
\textbf{N°} & \textbf{Comportement attendu} & \textbf{État} & \textbf{Remarques} \\
        \midrule
        \endfirsthead
        \caption[]{Bloc 4 -- Statistics / Storage (suite)} \\
        \toprule
        \textbf{N°} & \textbf{Comportement attendu} & \textbf{État} & \textbf{Remarques} \\
        \midrule
        \endhead
        \midrule
        \multicolumn{4}{r@{}}{\textit{Suite page suivante}} \\
        \endfoot
        \bottomrule
        \endlastfoot
        4.1 & Présence de la carte correctement indiquée, carte présente puis absente             & \OK & \\
        4.2 & Fréquence \gls{spi} affichée égale à 40 MHz                                         & \OK & \\
        4.3 & Capacité totale et utilisée conformes à celles vues sur un PC                       & \OK & \\
        4.4 & Nombre de pistes égal au nombre de fichiers \texttt{.wav} et \texttt{.mp3} présents sur la carte & \OK & \\
        4.5 & Le retrait de la carte met l'affichage à jour                                       & \OK & \\
        \end{longtable}
        }

\subsection{Bloc 5. Statistics / Battery}

Protocole : batterie branchée.

{\footnotesize
\setlength{\tabcolsep}{4pt}
\renewcommand{\arraystretch}{1.15}
\begin{longtable}{@{}
>{\raggedright\arraybackslash}p{0.045\textwidth}
>{\raggedright\arraybackslash}p{0.400\textwidth}
>{\centering\arraybackslash}p{0.055\textwidth}
>{\raggedright\arraybackslash}p{0.390\textwidth}@{}}
\caption{Bloc 5 -- Statistics / Battery}
\label{tab:ann_mesures_sw_b5} \\
\toprule
\textbf{N°} & \textbf{Comportement attendu} & \textbf{État} & \textbf{Remarques} \\
        \midrule
        \endfirsthead
        \caption[]{Bloc 5 -- Statistics / Battery (suite)} \\
        \toprule
        \textbf{N°} & \textbf{Comportement attendu} & \textbf{État} & \textbf{Remarques} \\
        \midrule
        \endhead
        \midrule
        \multicolumn{4}{r@{}}{\textit{Suite page suivante}} \\
        \endfoot
        \bottomrule
        \endlastfoot
        5.1 & L'état de charge affiché suit le branchement et le débranchement de l'USB & \OK & \\
        5.2 & \texttt{INOKB} au niveau haut, USB débranché                              & \OK & \\
        5.3 & Mux USB orienté vers le CP2102N, USB débranché                            & \OK & \\
        5.4 & \texttt{INOKB} passe au niveau bas au branchement de l'USB                & \OK & \\
        5.5 & Mux USB orienté vers le MAX77757 au branchement                           & \OK & \\
        5.6 & Mux USB revenu sur le CP2102N au plus tard 10 s après le branchement      & \OK & \\
        \end{longtable}
        }

        \subsection{Bloc 6. Statistics / Battery / Debug}

{\footnotesize
\setlength{\tabcolsep}{4pt}
\renewcommand{\arraystretch}{1.15}
\begin{longtable}{@{}
>{\raggedright\arraybackslash}p{0.045\textwidth}
>{\raggedright\arraybackslash}p{0.400\textwidth}
>{\centering\arraybackslash}p{0.055\textwidth}
>{\raggedright\arraybackslash}p{0.390\textwidth}@{}}
\caption{Bloc 6 -- Statistics / Battery / Debug}
\label{tab:ann_mesures_sw_b6} \\
\toprule
\textbf{N°} & \textbf{Comportement attendu} & \textbf{État} & \textbf{Remarques} \\
        \midrule
        \endfirsthead
        \caption[]{Bloc 6 -- Statistics / Battery / Debug (suite)} \\
        \toprule
        \textbf{N°} & \textbf{Comportement attendu} & \textbf{État} & \textbf{Remarques} \\
        \midrule
        \endhead
        \midrule
        \multicolumn{4}{r@{}}{\textit{Suite page suivante}} \\
        \endfoot
        \bottomrule
        \endlastfoot
        6.1 & Les registres de configuration de la jauge relus sont tous conformes                & \OK & \\
        6.2 & FullCapR compris entre 855 et 1045 mAh, soit la DesignCap de 950 mAh à $\pm$ 10 \%   & \OK & FullCapR = 1027 mAh. La valeur dépend du modèle ModelGauge configuré et non du seul circuit \cite{analog_MAX17260} \\
        \end{longtable}
        }

        \subsection{Bloc 7. Storage}

        Protocole : carte microSD de référence insérée.

{\footnotesize
\setlength{\tabcolsep}{4pt}
\renewcommand{\arraystretch}{1.15}
\begin{longtable}{@{}
>{\raggedright\arraybackslash}p{0.045\textwidth}
>{\raggedright\arraybackslash}p{0.400\textwidth}
>{\centering\arraybackslash}p{0.055\textwidth}
>{\raggedright\arraybackslash}p{0.390\textwidth}@{}}
\caption{Bloc 7 -- Storage}
\label{tab:ann_mesures_sw_b7} \\
\toprule
\textbf{N°} & \textbf{Comportement attendu} & \textbf{État} & \textbf{Remarques} \\
        \midrule
        \endfirsthead
        \caption[]{Bloc 7 -- Storage (suite)} \\
        \toprule
        \textbf{N°} & \textbf{Comportement attendu} & \textbf{État} & \textbf{Remarques} \\
        \midrule
        \endhead
        \midrule
        \multicolumn{4}{r@{}}{\textit{Suite page suivante}} \\
        \endfoot
        \bottomrule
        \endlastfoot
        7.1 & L'arborescence affichée correspond à celle vue sur un PC, sous-dossiers compris & \OK & \\
        7.2 & Tous les fichiers \texttt{.mp3} présents sont listés, et aucun fichier fantôme  & \OK & \\
        7.3 & Tous les fichiers \texttt{.wav} présents sont listés, et aucun fichier fantôme  & \OK & \\
        7.4 & Les fichiers d'une autre extension ne sont pas affichés                         & \OK & \\
        7.5 & La validation sur une piste ouvre la popup avec les deux actions                & \OK & \\
        7.6 & Play lance la lecture et bascule sur le menu Music                              & \OK & \\
        \end{longtable}
        }

        \subsection{Bloc 8. Music}

        Protocole : piste de référence, casque branché sur le jack, potentiomètre à mi-course.

{\footnotesize
\setlength{\tabcolsep}{4pt}
\renewcommand{\arraystretch}{1.15}
\begin{longtable}{@{}
>{\raggedright\arraybackslash}p{0.045\textwidth}
>{\raggedright\arraybackslash}p{0.400\textwidth}
>{\centering\arraybackslash}p{0.055\textwidth}
>{\raggedright\arraybackslash}p{0.390\textwidth}@{}}
\caption{Bloc 8 -- Music}
\label{tab:ann_mesures_sw_b8} \\
\toprule
\textbf{N°} & \textbf{Comportement attendu} & \textbf{État} & \textbf{Remarques} \\
        \midrule
        \endfirsthead
        \caption[]{Bloc 8 -- Music (suite)} \\
        \toprule
        \textbf{N°} & \textbf{Comportement attendu} & \textbf{État} & \textbf{Remarques} \\
        \midrule
        \endhead
        \midrule
        \multicolumn{4}{r@{}}{\textit{Suite page suivante}} \\
        \endfoot
        \bottomrule
        \endlastfoot
        8.1 & Le titre affiché correspond au fichier lancé                                                                                 & \OK & \\
        8.2 & Conteneur, profondeur, fréquence et débit affichés correspondent à ceux du fichier                                           & \OK & \\
        8.3 & Depuis la piste n, next passe à n+1 et prev à n-1. Comportement en début et en fin de playlist conforme à la spécification      & \OK & \\
        8.4 & Le son sort sur la sortie sélectionnée et sur elle seule                                                                     & \OK & \\
        8.5 & La barre de volume suit le potentiomètre sans saut ni oscillation                                                            & \OK & \\
        8.6 & Le registre de volume du \gls{DAC} varie de façon monotone entre les deux butées                                             & \OK & \\
        8.7 & En fin de piste, la lecture enchaîne sur la suivante sans coupure ni doublon                                                 & \OK & \\
        8.8 & Music / OutputSelection, l'icône change, le son cesse sur l'ancienne sortie et reprend sur la nouvelle, la position de lecture est conservée & \OK & \\
        \end{longtable}
        }

        \subsection{Bloc 9. Settings / Bluetooth}

        Protocole : appareil Bluetooth de référence à portée.

{\footnotesize
\setlength{\tabcolsep}{4pt}
\renewcommand{\arraystretch}{1.15}
\begin{longtable}{@{}
>{\raggedright\arraybackslash}p{0.045\textwidth}
>{\raggedright\arraybackslash}p{0.400\textwidth}
>{\centering\arraybackslash}p{0.055\textwidth}
>{\raggedright\arraybackslash}p{0.390\textwidth}@{}}
\caption{Bloc 9 -- Settings / Bluetooth}
\label{tab:ann_mesures_sw_b9} \\
\toprule
\textbf{N°} & \textbf{Comportement attendu} & \textbf{État} & \textbf{Remarques} \\
        \midrule
        \endfirsthead
        \caption[]{Bloc 9 -- Settings / Bluetooth (suite)} \\
        \toprule
        \textbf{N°} & \textbf{Comportement attendu} & \textbf{État} & \textbf{Remarques} \\
        \midrule
        \endhead
        \midrule
        \multicolumn{4}{r@{}}{\textit{Suite page suivante}} \\
        \endfoot
        \bottomrule
        \endlastfoot
        9.1 & Le scan liste les appareils à proximité, dont l'appareil de référence               & \OK & \\
        9.2 & L'appairage aboutit et l'état passe à connecté                                      & \OK & \\
        9.3 & La déconnexion volontaire aboutit et l'état est mis à jour                          & \OK & \\
        9.4 & La reconnexion à un appareil connu aboutit sans nouvel appairage                    & \OK & \\
        9.5 & L'oubli d'un appareil le retire de la liste, et un nouvel appairage est alors exigé & \OK & \\
        9.6 & L'audio sort sur l'appareil Bluetooth, les deux voies présentes et non inversées    & \OK & \\
        \end{longtable}
        }

        \subsection{Bloc 10. Settings / Audio}

{\footnotesize
\setlength{\tabcolsep}{4pt}
\renewcommand{\arraystretch}{1.15}
\begin{longtable}{@{}
>{\raggedright\arraybackslash}p{0.045\textwidth}
>{\raggedright\arraybackslash}p{0.400\textwidth}
>{\centering\arraybackslash}p{0.055\textwidth}
>{\raggedright\arraybackslash}p{0.390\textwidth}@{}}
\caption{Bloc 10 -- Settings / Audio}
\label{tab:ann_mesures_sw_b10} \\
\toprule
\textbf{N°} & \textbf{Comportement attendu} & \textbf{État} & \textbf{Remarques} \\
        \midrule
        \endfirsthead
        \caption[]{Bloc 10 -- Settings / Audio (suite)} \\
        \toprule
        \textbf{N°} & \textbf{Comportement attendu} & \textbf{État} & \textbf{Remarques} \\
        \midrule
        \endhead
        \midrule
        \multicolumn{4}{r@{}}{\textit{Suite page suivante}} \\
        \endfoot
        \bottomrule
        \endlastfoot
        10.1 & Bascule jack vers Bluetooth et retour, sans blocage ni redémarrage                & \OK & \\
        10.2 & Balance à fond à gauche, seule la voie gauche est audible, et réciproquement      & \OK & \\
        \end{longtable}
        }

        \subsection{Bloc 11. Settings / Screen}

        Protocole : USB débranché, pour que les temporisations d'extinction ne soient pas inhibées par la présence de l'alimentation externe.

{\footnotesize
\setlength{\tabcolsep}{4pt}
\renewcommand{\arraystretch}{1.15}
\begin{longtable}{@{}
>{\raggedright\arraybackslash}p{0.045\textwidth}
>{\raggedright\arraybackslash}p{0.400\textwidth}
>{\centering\arraybackslash}p{0.055\textwidth}
>{\raggedright\arraybackslash}p{0.390\textwidth}@{}}
\caption{Bloc 11 -- Settings / Screen}
\label{tab:ann_mesures_sw_b11} \\
\toprule
\textbf{N°} & \textbf{Comportement attendu} & \textbf{État} & \textbf{Remarques} \\
        \midrule
        \endfirsthead
        \caption[]{Bloc 11 -- Settings / Screen (suite)} \\
        \toprule
        \textbf{N°} & \textbf{Comportement attendu} & \textbf{État} & \textbf{Remarques} \\
        \midrule
        \endhead
        \midrule
        \multicolumn{4}{r@{}}{\textit{Suite page suivante}} \\
        \endfoot
        \bottomrule
        \endlastfoot
        11.1 & La luminosité varie visiblement entre les deux réglages extrêmes                                            & \OK & \\
        11.2 & Pour un réglage de n minutes, l'écran s'éteint dans la fenêtre attendue, testé sur deux réglages différents & \OK & \\
        \end{longtable}
        }

        \subsection{Bloc 12. Settings / Power}

        Protocole : USB débranché, pour que les temporisations d'extinction ne soient pas inhibées par la présence de l'alimentation externe.

{\footnotesize
\setlength{\tabcolsep}{4pt}
\renewcommand{\arraystretch}{1.15}
\begin{longtable}{@{}
>{\raggedright\arraybackslash}p{0.045\textwidth}
>{\raggedright\arraybackslash}p{0.400\textwidth}
>{\centering\arraybackslash}p{0.055\textwidth}
>{\raggedright\arraybackslash}p{0.390\textwidth}@{}}
\caption{Bloc 12 -- Settings / Power}
\label{tab:ann_mesures_sw_b12} \\
\toprule
\textbf{N°} & \textbf{Comportement attendu} & \textbf{État} & \textbf{Remarques} \\
        \midrule
        \endfirsthead
        \caption[]{Bloc 12 -- Settings / Power (suite)} \\
        \toprule
        \textbf{N°} & \textbf{Comportement attendu} & \textbf{État} & \textbf{Remarques} \\
        \midrule
        \endhead
        \midrule
        \multicolumn{4}{r@{}}{\textit{Suite page suivante}} \\
        \endfoot
        \bottomrule
        \endlastfoot
        12.1 & L'extinction automatique intervient dans la fenêtre du réglage, testée sur deux réglages différents & \OK & \\
        12.2 & L'extinction manuelle coupe le rail \texttt{+3.3V}                                                  & \OK & \\
        \end{longtable}
        }

        \subsection{Bloc 13. Cas dégradés}

{\footnotesize
\setlength{\tabcolsep}{4pt}
\renewcommand{\arraystretch}{1.15}
\begin{longtable}{@{}
>{\raggedright\arraybackslash}p{0.045\textwidth}
>{\raggedright\arraybackslash}p{0.400\textwidth}
>{\centering\arraybackslash}p{0.055\textwidth}
>{\raggedright\arraybackslash}p{0.390\textwidth}@{}}
\caption{Bloc 13 -- Cas dégradés}
\label{tab:ann_mesures_sw_b13} \\
\toprule
\textbf{N°} & \textbf{Comportement attendu} & \textbf{État} & \textbf{Remarques} \\
        \midrule
        \endfirsthead
        \caption[]{Bloc 13 -- Cas dégradés (suite)} \\
        \toprule
        \textbf{N°} & \textbf{Comportement attendu} & \textbf{État} & \textbf{Remarques} \\
        \midrule
        \endhead
        \midrule
        \multicolumn{4}{r@{}}{\textit{Suite page suivante}} \\
        \endfoot
        \bottomrule
        \endlastfoot
        13.1 & Retrait de la carte SD pendant la lecture, arrêt propre et pas de redémarrage                                       & \OK  & \\
        13.2 & Extinction de l'enceinte Bluetooth pendant la lecture, perte de lien détectée et retour à un état cohérent          & \OK  & \\
        13.3 & Fichier \texttt{.wav} tronqué ou de taille nulle, refus ou arrêt propre, pas de blocage                            & \NOK & Le firmware lance la lecture au lieu de signaler une erreur. La pipeline audio reste fonctionnelle \\
        13.4 & Fichier d'extension supportée mais de format non géré, message explicite et pas de blocage                         & \NOK & Le firmware lance la lecture au lieu de signaler une erreur. La pipeline audio reste fonctionnelle \\
        13.5 & Démarrage sans carte SD, menu Storage cohérent et pas de blocage                                                    & \OK  & \\
        13.6 & Fichier au nom long et accentué, affichage correct et lecture possible                                             & \OK  & \\
        13.7 & Appuis rapides et répétés sur les boutons, aucun événement perdu et pas de saturation de la file                    & \OK  & \\
        13.8 & Décharge sous le seuil de sécurité, extinction propre et pas de corruption de la carte                             & \OK  & \\
        \end{longtable}
        }

        \subsection{Bloc 14. Robustesse et durée}

{\footnotesize
\setlength{\tabcolsep}{4pt}
\renewcommand{\arraystretch}{1.15}
\begin{longtable}{@{}
>{\raggedright\arraybackslash}p{0.045\textwidth}
>{\raggedright\arraybackslash}p{0.400\textwidth}
>{\centering\arraybackslash}p{0.055\textwidth}
>{\raggedright\arraybackslash}p{0.390\textwidth}@{}}
\caption{Bloc 14 -- Robustesse et durée}
\label{tab:ann_mesures_sw_b14} \\
\toprule
\textbf{N°} & \textbf{Comportement attendu} & \textbf{État} & \textbf{Remarques} \\
        \midrule
        \endfirsthead
        \caption[]{Bloc 14 -- Robustesse et durée (suite)} \\
        \toprule
        \textbf{N°} & \textbf{Comportement attendu} & \textbf{État} & \textbf{Remarques} \\
        \midrule
        \endhead
        \midrule
        \multicolumn{4}{r@{}}{\textit{Suite page suivante}} \\
        \endfoot
        \bottomrule
        \endlastfoot
        14.1 & Appui sur S1 pendant l'affichage, redémarrage propre, flash et \gls{psram} opérationnelles                & \OK & \\
        14.2 & Coupure d'alimentation brutale en lecture, au redémarrage le marqueur \gls{nvs} et le verdict d'extinction sont cohérents & \OK & \\
        14.3 & Extinction normale, verdict d'extinction nominal au redémarrage suivant                                   & \OK & \\
        \end{longtable}
        }

        \subsection{Bloc 15. Absence de bruits parasites}

        Protocole : évaluation à l'oreille, casque branché sur le jack.

{\footnotesize
\setlength{\tabcolsep}{4pt}
\renewcommand{\arraystretch}{1.15}
\begin{longtable}{@{}
>{\raggedright\arraybackslash}p{0.045\textwidth}
>{\raggedright\arraybackslash}p{0.400\textwidth}
>{\centering\arraybackslash}p{0.055\textwidth}
>{\raggedright\arraybackslash}p{0.390\textwidth}@{}}
\caption{Bloc 15 -- Absence de bruits parasites}
\label{tab:ann_mesures_sw_b15} \\
\toprule
\textbf{N°} & \textbf{Comportement attendu} & \textbf{État} & \textbf{Remarques} \\
        \midrule
        \endfirsthead
        \caption[]{Bloc 15 -- Absence de bruits parasites (suite)} \\
        \toprule
        \textbf{N°} & \textbf{Comportement attendu} & \textbf{État} & \textbf{Remarques} \\
        \midrule
        \endhead
        \midrule
        \multicolumn{4}{r@{}}{\textit{Suite page suivante}} \\
        \endfoot
        \bottomrule
        \endlastfoot
        15.1 & Aucun pop audible à l'allumage                          & \OK  & \\
        15.2 & Aucun pop audible à l'extinction                        & \OK  & \\
        15.3 & Aucun pop audible au début et à la fin d'une lecture    & \OK  & \\
        15.4 & Aucun pop audible au changement de sortie                & \OK  & \\
        15.5 & Aucun pop audible au changement de piste                 & \NOK & Un léger glitch reste audible. L'origine, logicielle ou matérielle, n'a pas pu être établie faute de temps \\
    \end{longtable}
}

\section{Matrice des formats audio}
\label{ann:mesures_formats}

La piste de référence est générée par script Python, à l'aide de \texttt{numpy} et de \texttt{soundfile}, et déclinée à l'identique dans tous les formats. Elle enchaîne 10 s de 1 kHz sur les deux voies, 10 s sur la voie gauche seule, 10 s sur la voie droite seule, 2 s de silence, puis un balayage jusqu'à la fin du fichier.

Cette piste unique valide d'un seul essai le mapping des voies, l'absence de coupure et la lecture jusqu'à la fin. L'écart entre la durée annoncée par le conteneur et la durée chronométrée révèle une erreur d'horloge, le critère retenu étant un écart inférieur à 1 \%.

\begin{table}[H]
    \centering
    \caption{Formats lus correctement sur chacune des deux sorties}
    \label{tab:ann_formats}
    \small
    \setlength{\tabcolsep}{5pt}
    \begin{tabular}{@{}l l l r c c r r r@{}}
        \toprule
                    &                    &                     &                & \multicolumn{2}{c}{\textbf{Sortie}} & \multicolumn{3}{c}{\textbf{Durée [s]}}                                                         \\
        \cmidrule(lr){5-6} \cmidrule(lr){7-9}
        \textbf{N°} & \textbf{Conteneur} & \textbf{Profondeur} & \textbf{Fréq.} & \textbf{Jack}                       & \textbf{BT}                            & \textbf{Mesurée} & \textbf{Affichée} & \textbf{Écart} \\
        \midrule
        F1          & WAV                & 16 bits             & 44,1 kHz       & \OK                                 & \OK                                    & 42               & 42                & 0              \\
        F2          & WAV                & 16 bits             & 48 kHz         & \OK                                 & --                                     & 42               & 42                & 0              \\
        F3          & WAV                & 16 bits             & 88,2 kHz       & \OK                                 & --                                     & 42               & 42                & 0              \\
        F4          & WAV                & 16 bits             & 96 kHz         & \OK                                 & --                                     & 42               & 42                & 0              \\
        F5          & WAV                & 16 bits             & 176,4 kHz      & \OK                                 & --                                     & 42               & 42                & 0              \\
        F6          & WAV                & 16 bits             & 192 kHz        & \OK                                 & --                                     & 42               & 42                & 0              \\
        F7          & WAV                & 24 bits             & 44,1 kHz       & \OK                                 & \OK                                    & 42               & 42                & 0              \\
        F8          & WAV                & 24 bits             & 48 kHz         & \OK                                 & --                                     & 42               & 42                & 0              \\
        F9          & WAV                & 24 bits             & 88,2 kHz       & \OK                                 & --                                     & 42               & 42                & 0              \\
        F10         & WAV                & 24 bits             & 96 kHz         & \OK                                 & --                                     & 42               & 42                & 0              \\
        F11         & WAV                & 24 bits             & 176,4 kHz      & \OK                                 & --                                     & 42               & 42                & 0              \\
        F12         & WAV                & 24 bits             & 192 kHz        & \OK                                 & --                                     & 42               & 42                & 0              \\
        F13         & MP3                & 128 kbit/s          & 44,1 kHz       & \OK                                 & \OK                                    & 42               & 42                & 0              \\
        F14         & MP3                & 192 kbit/s          & 44,1 kHz       & \OK                                 & \OK                                    & 42               & 42                & 0              \\
        F15         & MP3                & 320 kbit/s          & 44,1 kHz       & \OK                                 & \OK                                    & 42               & 42                & 0              \\
        F16         & MP3                & 320 kbit/s          & 48 kHz         & \OK                                 & --                                     & 42               & 42                & 0              \\
        \bottomrule
    \end{tabular}
\end{table}

La sortie jack accepte les seize formats. La sortie Bluetooth se limite pour l'instant aux fichiers à 44,1 kHz, fréquence unique gérée par la chaîne \gls{a2dp} telle qu'elle est implémentée. Un tiret indique donc une combinaison non prise en charge, non un essai en échec.

\section{Caractérisation de la sortie filaire}
\label{ann:mesures_audio}

La chaîne mesurée est celle du schéma, PCM5242 en \gls{i2s} 3 fils sans \gls{mclk}, filtre RC différentiel, MAX97220A en gain 0,5 et jack 3,5 mm. La pleine échelle du \gls{DAC} vaut 4,2 Vrms différentiels en mode VREF \cite{TI_PCM5242}, soit 2,1 Vrms au jack à atténuation numérique nulle.

Les blocs A et B sont relevés au multimètre et à l'oscilloscope. Les blocs C à F le sont à l'Audio Precision ATS-2. Toute la campagne se déroule dans une configuration unique, volume système au maximum et registre de volume du \gls{DAC} à 0 dB, le niveau de sollicitation étant fixé par le seul choix du fichier joué. Les relevés chiffrés associés figurent à la section~\ref{ann:mesures_ap}.

\subsection{Jeux de fichiers de test}

Tous les fichiers sont produits par le script \texttt{gen\_fichiers\_test\_ap.py}. Le tableau~\ref{tab:ann_jeux_fichiers} donne la composition de chaque jeu et sa convention de nom, ce qui permet de désigner les fichiers par leur seul jeu et leur niveau dans les tableaux de relevés qui suivent.

\begin{table}[H]
    \centering
    \caption{Jeux de fichiers de test et conventions de nom}
    \label{tab:ann_jeux_fichiers}
    \small
    \setlength{\tabcolsep}{5pt}
    \begin{tabular}{@{}
        >{\raggedright\arraybackslash}p{0.10\textwidth}
        >{\raggedright\arraybackslash}p{0.42\textwidth}
        >{\raggedright\arraybackslash}p{0.34\textwidth}@{}}
        \toprule
        \textbf{Jeu} & \textbf{Contenu}                                                & \textbf{Convention de nom}               \\
        \midrule
        THDN         & 1 kHz aux niveaux 0, $-1$ et $-20$ \gls{dbfs}                   & \texttt{THDN\_1000Hz\_...dBFS.wav}       \\
        FR\_44k      & 13 fréquences de 20 Hz à 20 kHz, source à 44,1 kHz et $-1$ dBFS & \texttt{FR\_44k\_...Hz\_-1dBFS.wav}      \\
        LEVEL        & 1 kHz, 8 niveaux de 0 à $-60$ dBFS                              & \texttt{LEVEL\_1000Hz\_...dBFS.wav}      \\
        XT\_L, XT\_R & 1 kHz et 10 kHz à 0 dBFS sur une seule voie                     & \texttt{XT\_L\_...Hz\_0dBFS.wav}         \\
        MP3          & Référence WAV et versions MP3 à 128 et 320 kbit/s               & \texttt{REF.wav}, \texttt{MP3\_320k.mp3} \\
        Silence      & Silence numérique, 44,1 kHz et 24 bits                          & \texttt{Silence\_44k\_24bit.wav}         \\
        filtre       & Balayages de 1 à 90 kHz, source à 192 kHz et $-6$ dBFS          & \texttt{SWEEP\_...kHz\_192k\_-6dBFS.wav} \\
        \bottomrule
    \end{tabular}
\end{table}

\subsection{Bloc A. Vérification du gain de la chaîne analogique}

Protocole : fichier THDN à $-20$ dBFS lu depuis la microSD. Volume système au maximum et registre de volume du \gls{DAC} à 0 dB. La marge nécessaire vient du seul niveau du fichier, ce qui laisse 21 dB sous les rails et rend tout écrêtage impossible. Sortie jack sélectionnée, sortie à vide. Sonde 1 sur \texttt{OUTLP} et sonde 2 sur \texttt{OUTLN}, mesure différentielle A moins B pour la sortie du \gls{DAC}. Sonde sur la pointe du jack pour la sortie de l'amplificateur.

{\scriptsize
\setlength{\tabcolsep}{3pt}
\renewcommand{\arraystretch}{1.15}
\begin{longtable}{@{}
    >{\raggedright\arraybackslash}p{0.042\textwidth}
    >{\raggedright\arraybackslash}p{0.220\textwidth}
    >{\raggedright\arraybackslash}p{0.180\textwidth}
    >{\raggedright\arraybackslash}p{0.115\textwidth}
    >{\centering\arraybackslash}p{0.052\textwidth}
    >{\raggedright\arraybackslash}p{0.268\textwidth}@{}}
    \caption{Bloc A -- Vérification du gain de la chaîne analogique}
    \label{tab:ann_mesures_audio_bA}                                                                                                                                                                                                  \\
    \toprule
    \textbf{N°} & \textbf{Point de mesure}                                                            & \textbf{Valeur attendue} & \textbf{Mesure} & \textbf{État} & \textbf{Remarques}                                               \\
    \midrule
    \endfirsthead
    \caption[]{Bloc A -- Vérification du gain de la chaîne analogique (suite)}                                                                                                                                                        \\
    \toprule
    \textbf{N°} & \textbf{Point de mesure}                                                            & \textbf{Valeur attendue} & \textbf{Mesure} & \textbf{État} & \textbf{Remarques}                                               \\
    \midrule
    \endhead
    \midrule
    \multicolumn{6}{r@{}}{\textit{Suite page suivante}}                                                                                                                                                                               \\
    \endfoot
    \bottomrule
    \endlastfoot
    A.1         & Amplitude différentielle en sortie du \gls{DAC}, voie gauche, jeu THDN à $-20$ dBFS & 420 mVrms                & 422 mVrms       & \OK           & Pleine échelle de 4,2 Vrms réduite par le seul niveau du fichier \\
    A.2         & Amplitude en sortie du jack, voie gauche                                            & 210 mVrms                & 211 mVrms       & \OK           &                                                                  \\
    A.3         & Gain de la voie gauche, rapport de A.2 sur A.1                                      & 0,5 $\pm$ 5 \%           & 0,500           & \OK           & Conforme au dimensionnement                                      \\
    A.4         & Amplitude différentielle en sortie du \gls{DAC}, voie droite                        & 420 mVrms                & 422 mVrms       & \OK           &                                                                  \\
    A.5         & Amplitude en sortie du jack, voie droite                                            & 210 mVrms                & 212 mVrms       & \OK           &                                                                  \\
    A.6         & Gain de la voie droite, rapport de A.5 sur A.4                                      & 0,5 $\pm$ 5 \%           & 0,502           & \OK           &                                                                  \\
\end{longtable}
}

\begin{figure}[H]
    \centering
    \begin{subfigure}{0.49\textwidth}
        \centering
        \includegraphics[width=\textwidth]{\assetsdir/sortie_diff_dac_left.png}
        \caption{Sortie différentielle du \gls{DAC}, voie gauche, 12 Vpp}
    \end{subfigure}
    \hfill
    \begin{subfigure}{0.49\textwidth}
        \centering
        \includegraphics[width=\textwidth]{\assetsdir/sortie_jack_left.png}
        \caption{Sortie du jack, voie gauche, 6 Vpp}
    \end{subfigure}
    \caption{Amplitudes relevées de part et d'autre de l'amplificateur, fichier à 1 kHz et 0 dBFS, sortie à vide. Le rapport des deux amplitudes vaut 0,5, conforme au dimensionnement, et la sinusoïde reste propre sur les deux captures, sans aplatissement des crêtes ni inversion de phase.}
    \label{fig:gain_chaine_analogique}
\end{figure}

Deux captures complémentaires, prises hors de la campagne numérotée, illustrent le comportement de l'étage analogique à pleine échelle et sortie à vide.

La figure~\ref{fig:gain_chaine_analogique} place côte à côte la sortie différentielle du \gls{DAC} et la sortie du jack pour le même fichier à 1 kHz et 0 dBFS. On voit clairement contraste avec la figure~\ref{saturation_negative.PNG}, où la même chaîne, au même niveau numérique, déforme le signal lorsqu'un casque de 32 $\Omega$ est branché.

\subsection{Bloc B. Nature de l'écrêtage à pleine échelle}

Protocole : mêmes réglages qu'au bloc A. Le bloc se parcourt en deux passes qui ne diffèrent que par le fichier joué, d'abord THDN à $-20$ dBFS qui ne peut pas écrêter, puis THDN à 0 dBFS qui sollicite la pleine échelle de conception. Sortie à vide, puis chargée par un casque de 32 $\Omega$.

{\scriptsize
    \setlength{\tabcolsep}{3pt}
    \renewcommand{\arraystretch}{1.15}
    \begin{longtable}{@{}
        >{\raggedright\arraybackslash}p{0.042\textwidth}
        >{\raggedright\arraybackslash}p{0.220\textwidth}
        >{\raggedright\arraybackslash}p{0.180\textwidth}
        >{\raggedright\arraybackslash}p{0.115\textwidth}
        >{\centering\arraybackslash}p{0.052\textwidth}
        >{\raggedright\arraybackslash}p{0.268\textwidth}@{}}
        \caption{Bloc B -- Nature de l'écrêtage à pleine échelle}
        \label{tab:ann_mesures_audio_bB} \\
        \toprule
        \textbf{N°} & \textbf{Point de mesure} & \textbf{Valeur attendue} & \textbf{Mesure} & \textbf{État} & \textbf{Remarques} \\
        \midrule
        \endfirsthead
        \caption[]{Bloc B -- Nature de l'écrêtage à pleine échelle (suite)} \\
        \toprule
        \textbf{N°} & \textbf{Point de mesure} & \textbf{Valeur attendue} & \textbf{Mesure} & \textbf{État} & \textbf{Remarques} \\
        \midrule
        \endhead
        \midrule
        \multicolumn{6}{r@{}}{\textit{Suite page suivante}} \\
        \endfoot
        \bottomrule
        \endlastfoot
        B.1  & Amplitude en sortie du jack, jeu THDN à $-20$ dBFS      & 210 mVrms                                                                              & L 211 mVrms\newline R 212 mVrms                                & \OK  & \\
        B.2  & Forme d'onde en sortie du jack, $-20$ dBFS               & Sinusoïde propre, aucun aplatissement                                                 & Sinusoïde propre                                               & \OK  & \\
        B.3  & Forme d'onde différentielle en sortie du \gls{DAC}, $-20$ dBFS & Sinusoïde propre                                                                & Sinusoïde propre, paliers du \gls{DAC} légèrement perceptibles & \OK  & \\
        B.4  & Amplitude en sortie du jack, jeu THDN à 0 dBFS, à vide   & 2,1 Vrms                                                                              & L 2,13 Vrms\newline R 2,13 Vrms                                & \OK  & \\
        B.5  & Forme d'onde en sortie du jack, 0 dBFS, à vide           & Sinusoïde propre                                                                      & Sinusoïde propre                                               & \OK  & \\
        B.6  & Forme d'onde différentielle en sortie du \gls{DAC}, 0 dBFS & Sinusoïde propre                                                                    & Sinusoïde propre, paliers du \gls{DAC} légèrement perceptibles & \OK  & \\
        B.7  & Amplitude en sortie du jack, 0 dBFS, casque de 32 $\Omega$ & 2,1 Vrms                                                                            & 1,83 Vrms                                                      & \NOK & Le niveau s'effondre dès que la sortie est chargée \\
        B.8  & Forme d'onde en sortie du jack, 0 dBFS, casque de 32 $\Omega$ & Sinusoïde propre                                                                 & Écrêtage de l'alternance négative                              & \NOK & L'alternance positive reste lisse, l'alternance négative présente un plat marqué, figure~\ref{saturation_negative.png}. La cause est cherchée au point B.10 \\
        B.9  & PVSS du MAX97220A, pin 5, 0 dBFS, à vide                 & $-3,3$ V, identique au point 10.8 du tableau~\ref{tab:ann_mesures_hw_b10}             & $-3,291$ V                                                     & \OK  & \\
        B.10 & PVSS du MAX97220A, pin 5, 0 dBFS, casque de 32 $\Omega$  & $-3,3$ V                                                                              & $-2,720$ V                                                     & \NOK & Mesure décisive du bloc. La pompe de charge ne tient pas le rail négatif sous charge \cite{analog_MAX97220A}, ce qui explique que l'écrêtage observé au point B.8 n'affecte que l'alternance négative \\
        B.11 & Niveau numérique maximal sans écrêtage, casque de 32 $\Omega$, jeu LEVEL & Premier niveau, en descendant depuis 0 dBFS, où l'alternance négative redevient lisse & $-6$ dBFS                                          & \OK  & Conduit à limiter le volume maximal du lecteur à $-6$ dB par le registre du \gls{DAC} \\
        \end{longtable}
        }

        Les points B.7, B.8 et B.10 forment le seul résultat vraiment structurant de cette section. La figure~\ref{saturation_negative.png} montre la signature caractéristique observée au point B.8, une alternance positive intacte et une alternance négative aplatie, dissymétrie qui oriente immédiatement le diagnostic vers le rail négatif et non vers l'étage numérique.

        \fig[H, width=0.9\textwidth]{Écrêtage de l'alternance négative en sortie du jack, fichier à 0 dBFS débité dans un casque de 32 $\Omega$}{saturation_negative.png}

        Cette dissymétrie exclut à elle seule deux causes possibles. Un écrêtage numérique en amont ou une saturation de l'étage différentiel du \gls{DAC} affecteraient les deux alternances de la même manière. Seule une faiblesse d'un des deux rails de l'amplificateur peut produire un aplatissement d'un seul côté, ce que le point B.10 confirme en mesurant l'effondrement de PVSS de $-3,291$ V à vide à $-2,720$ V sous charge.

        \subsection{Bloc C. Préparation et câblage de la session Audio Precision}

        Protocole : à dérouler avant toute mesure à l'analyseur. Ces points ne produisent aucun résultat exploitable mais conditionnent la validité de tous les suivants. Logiciel ATS 1.60, setups du dossier AP\_Setup.

{\scriptsize
\setlength{\tabcolsep}{3pt}
\renewcommand{\arraystretch}{1.15}
\begin{longtable}{@{}
>{\raggedright\arraybackslash}p{0.042\textwidth}
>{\raggedright\arraybackslash}p{0.220\textwidth}
>{\raggedright\arraybackslash}p{0.180\textwidth}
>{\raggedright\arraybackslash}p{0.115\textwidth}
>{\centering\arraybackslash}p{0.052\textwidth}
>{\raggedright\arraybackslash}p{0.268\textwidth}@{}}
\caption{Bloc C -- Préparation et câblage de la session Audio Precision}
\label{tab:ann_mesures_audio_bC} \\
\toprule
\textbf{N°} & \textbf{Point de mesure} & \textbf{Valeur attendue} & \textbf{Mesure} & \textbf{État} & \textbf{Remarques} \\
\midrule
\endfirsthead
\caption[]{Bloc C -- Préparation et câblage de la session Audio Precision (suite)} \\
\toprule
\textbf{N°} & \textbf{Point de mesure} & \textbf{Valeur attendue} & \textbf{Mesure} & \textbf{État} & \textbf{Remarques} \\
        \midrule
        \endhead
        \midrule
        \multicolumn{6}{r@{}}{\textit{Suite page suivante}} \\
        \endfoot
        \bottomrule
        \endlastfoot
        C.1 & Câble jack 3,5 mm vers deux BNC, pointe sur Channel A et bague sur Channel B & Continuité vérifiée                                                          & Continuité confirmée sur les deux voies                        & \OK & Contrôlé à l'ohmmètre avant tout branchement de l'analyseur \\
        C.2 & Setup des mesures de niveau et de THD+N                                     & Le fichier \texttt{AP\_THD.ats2} s'ouvre                                     & Ouverture correcte                                              & \OK & \\
        C.3 & Setup des mesures de bruit et de spectre                                     & Le fichier \texttt{AP\_NOISE.ats2} s'ouvre                                   & Ouverture correcte                                              & \OK & \\
        C.4 & Contrôle de l'appariement des voies, jeux XT\_L et XT\_R                     & \texttt{XT\_L} joue sur la voie A seule, \texttt{XT\_R} sur la voie B seule  & Confirmé par le niveau lu sur chaque voie                       & \OK & Détecte une inversion de voies avant de commencer. Sans ce contrôle, les mesures de diaphonie du bloc E seraient ininterprétables \\
        C.5 & Volume système au maximum et registre de volume du \gls{DAC} à 0 dB          & Confirmé par la position du potentiomètre                                    & Configuration confirmée                                         & \OK & Toutes les valeurs attendues des blocs D et E supposent cette configuration, inchangée d'un bloc à l'autre \\
        C.6 & Setup des mesures de diaphonie                                                & Le fichier \texttt{AP\_XTALK.ats2} s'ouvre                                   & Ouverture correcte                                              & \OK & \\
        \end{longtable}
        }

        \subsection{Bloc D. Analyseur en configuration THD+N}

        Protocole : instrument Audio Analyzer, setup \texttt{AP\_THD.ats2}, entrée Analog, Measurement Function sur THD+N Ratio affichée en dB. Détection \gls{rms}, bande de 22 Hz en passe-haut et 20 kHz en passe-bas, filtre pondérateur désactivé. Entrées Channel A et B en Unbalanced, couplage continu, autorange. Les entrées de l'analyseur présentent 100 k$\Omega$, toutes ces mesures sont donc non chargées. La \gls{diaphonie} ne se mesure pas avec ces réglages, elle fait l'objet du bloc E et de son propre setup.

{\scriptsize
\setlength{\tabcolsep}{3pt}
\renewcommand{\arraystretch}{1.15}
\begin{longtable}{@{}
>{\raggedright\arraybackslash}p{0.042\textwidth}
>{\raggedright\arraybackslash}p{0.220\textwidth}
>{\raggedright\arraybackslash}p{0.180\textwidth}
>{\raggedright\arraybackslash}p{0.115\textwidth}
>{\centering\arraybackslash}p{0.052\textwidth}
>{\raggedright\arraybackslash}p{0.268\textwidth}@{}}
\caption{Bloc D -- Analyseur en configuration THD+N}
\label{tab:ann_mesures_audio_bD} \\
\toprule
\textbf{N°} & \textbf{Point de mesure} & \textbf{Valeur attendue} & \textbf{Mesure} & \textbf{État} & \textbf{Remarques} \\
\midrule
\endfirsthead
\caption[]{Bloc D -- Analyseur en configuration THD+N (suite)} \\
\toprule
\textbf{N°} & \textbf{Point de mesure} & \textbf{Valeur attendue} & \textbf{Mesure} & \textbf{État} & \textbf{Remarques} \\
        \midrule
        \endhead
        \midrule
        \multicolumn{6}{r@{}}{\textit{Suite page suivante}} \\
        \endfoot
        \bottomrule
        \endlastfoot
        D.1  & Niveau de sortie, 1 kHz à 0 dBFS                     & 2,1 Vrms                        & L 2,139 Vrms\newline R 2,143 Vrms & \OK & Pleine échelle de 4,2 Vrms différentiels et gain 0,5, soit un gain réel de 0,509 \\
        D.2  & Offset continu mesuré par l'analyseur                & $<$ 350 $\mu$V                  & L 17,5 $\mu$V\newline R 14,2 $\mu$V & \OK & Valeur plus basse qu'aux points 10.11 et 10.12 du tableau~\ref{tab:ann_mesures_hw_b10}, l'analyseur étant plus résolu que le multimètre \\
        D.3  & THD+N à 1 kHz, $-1$ dBFS                             & Environ $-94$ dB                & L $-94,7$ dB\newline R $-94,9$ dB & \OK & PCM5242 annoncé à $-94$ dB typique et $-87$ dB minimum \cite{TI_PCM5242}. Résidu propre de l'analyseur à $-101$ dB, donc non limitant \\
        D.4  & THD+N à 1 kHz, 0 dBFS                                & Non contraint                   & $-95$ dB                          & \OK & \\
        D.5  & Réponse en fréquence, 20 Hz à 20 kHz, jeu FR\_44k    & Écart inférieur à $\pm$ 0,1 dB  & Inférieur à $\pm$ 0,1 dB          & \OK & Série complète au tableau~\ref{tab:ann_ap_reponse} \\
        D.6  & THD+N en fonction de la fréquence, 20 Hz à 20 kHz    & Meilleure que $-90$ dB          & $-92,6$ dB                        & \OK & Pire valeur de la série, relevée à 5 kHz. Voir le tableau~\ref{tab:ann_ap_reponse} \\
        D.7  & THD+N en fonction du niveau, 1 kHz, jeu LEVEL        & Courbe à conserver              & Voir le tableau~\ref{tab:ann_ap_niveau} & \OK & Montre le croisement entre plancher de bruit et plancher de distorsion \\
        D.8  & Niveau relevé à 20 kHz, source à 44,1 kHz            & Fléchissement à relever         & $-0,014$ dB                       & \OK & Aucun fléchissement significatif en haut de bande \\
        D.9a & THD+N à 1 kHz, source WAV, jeu MP3                   & Proche du point D.3             & L $-96,3$ dB\newline R $-96,1$ dB & \OK & Ces mesures ont été faites un autre jour, l'écart avec D.3 reste dans l'incertitude de mesure \\
        D.9b & THD+N à 1 kHz, source MP3 à 320 kbit/s               & Dégradation par rapport à D.9a  & L $-87,8$ dB\newline R $-87,6$ dB & \OK & Le codage MP3 dégrade le THD+N d'environ 8 dB \\
        D.9c & THD+N à 1 kHz, source MP3 à 128 kbit/s               & Dégradation par rapport à D.9b  & L $-87,2$ dB\newline R $-87,1$ dB & \OK & L'écart entre 320 et 128 kbit/s reste faible sur ce critère \\
        \end{longtable}
        }

\subsection{Bloc E. Analyseur en configuration diaphonie}

Protocole : instrument Audio Analyzer, setup \texttt{AP\_XTALK.ats2}. La diaphonie n'est pas une mesure de distorsion mais un rapport de niveaux, à la fréquence d'essai, entre le niveau de la voie qui se tait et celui de la voie qui joue. La Measurement Function est réglée sur Crosstalk en dB, l'analyseur affichant alors directement ce rapport. Chaque fréquence est mesurée dans les deux sens. Le jeu \texttt{XT\_L} excite la voie gauche et la fuite se lit sur Channel B, le jeu \texttt{XT\_R} excite la voie droite et la fuite se lit sur Channel A. La valeur affichée sur la voie excitée n'a aucune signification physique et n'est pas consignée. Le câble de la voie muette est débranché côté carte.

{\scriptsize
\setlength{\tabcolsep}{3pt}
\renewcommand{\arraystretch}{1.15}
\begin{longtable}{@{}
>{\raggedright\arraybackslash}p{0.042\textwidth}
>{\raggedright\arraybackslash}p{0.220\textwidth}
>{\raggedright\arraybackslash}p{0.180\textwidth}
>{\raggedright\arraybackslash}p{0.115\textwidth}
>{\centering\arraybackslash}p{0.052\textwidth}
>{\raggedright\arraybackslash}p{0.268\textwidth}@{}}
\caption{Bloc E -- Analyseur en configuration diaphonie}
\label{tab:ann_mesures_audio_bE} \\
\toprule
\textbf{N°} & \textbf{Point de mesure} & \textbf{Valeur attendue} & \textbf{Mesure} & \textbf{État} & \textbf{Remarques} \\
        \midrule
        \endfirsthead
        \caption[]{Bloc E -- Analyseur en configuration diaphonie (suite)} \\
        \toprule
        \textbf{N°} & \textbf{Point de mesure} & \textbf{Valeur attendue} & \textbf{Mesure} & \textbf{État} & \textbf{Remarques} \\
        \midrule
        \endhead
        \midrule
        \multicolumn{6}{r@{}}{\textit{Suite page suivante}} \\
        \endfoot
        \bottomrule
        \endlastfoot
        E.1 & Niveau sur la voie excitée, contrôle de la référence, jeu XT\_L à 1 kHz & 2,1 Vrms                                & 2,139 Vrms   & \OK & Lu sur Channel A, la voie qui joue. Confirme que le fichier sort au niveau attendu \\
        E.2 & Diaphonie de la voie gauche vers la voie droite, 1 kHz                  & Environ $-90$ dB                        & $-102,6$ dB  & \OK & \\
        E.3 & Diaphonie de la voie droite vers la voie gauche, 1 kHz                  & Environ $-90$ dB                        & $-106,7$ dB  & \OK & \\
        E.4 & Diaphonie de la voie gauche vers la voie droite, 10 kHz                 & Dégradation attendue par rapport à E.2  & $-92,5$ dB   & \OK & Le couplage capacitif croît avec la fréquence \\
        E.5 & Diaphonie de la voie droite vers la voie gauche, 10 kHz                 & Dégradation attendue par rapport à E.3  & $-88,0$ dB   & \OK & \\
    \end{longtable}
}

\subsection{Bloc F. Analyseur en configuration bruit}

Protocole : instrument Audio Analyzer, setup \texttt{AP\_NOISE.ats2}, Measurement Function sur Amplitude. La détection \gls{rms} est impérative, le facteur de crête du bruit n'étant pas celui d'une sinusoïde. Le panneau Sweep sert aux analyses spectrales. Mesures non chargées sur les 100 k$\Omega$ de l'analyseur.

{\scriptsize
\setlength{\tabcolsep}{3pt}
\renewcommand{\arraystretch}{1.15}
\begin{longtable}{@{}
    >{\raggedright\arraybackslash}p{0.042\textwidth}
    >{\raggedright\arraybackslash}p{0.220\textwidth}
    >{\raggedright\arraybackslash}p{0.180\textwidth}
    >{\raggedright\arraybackslash}p{0.115\textwidth}
    >{\centering\arraybackslash}p{0.052\textwidth}
    >{\raggedright\arraybackslash}p{0.268\textwidth}@{}}
    \caption{Bloc F -- Analyseur en configuration bruit}
    \label{tab:ann_mesures_audio_bF}                                                                                                                                                                                                       \\
    \toprule
    \textbf{N°} & \textbf{Point de mesure}                                               & \textbf{Valeur attendue} & \textbf{Mesure}               & \textbf{État} & \textbf{Remarques}                                                   \\
    \midrule
    \endfirsthead
    \caption[]{Bloc F -- Analyseur en configuration bruit (suite)}                                                                                                                                                                         \\
    \toprule
    \textbf{N°} & \textbf{Point de mesure}                                               & \textbf{Valeur attendue} & \textbf{Mesure}               & \textbf{État} & \textbf{Remarques}                                                   \\
    \midrule
    \endhead
    \midrule
    \multicolumn{6}{r@{}}{\textit{Suite page suivante}}                                                                                                                                                                                    \\
    \endfoot
    \bottomrule
    \endlastfoot
    F.1         & Plancher de bruit pondéré A, jeu Silence                               & Relevé, en $\mu$V        & 5,252 $\mu$V                  & \OK           & Filtre sur A weighting, bande de 22 Hz à 20 kHz, détection \gls{rms} \\
    F.2         & Plancher de bruit non pondéré, 22 Hz à 20 kHz                          & Relevé, en $\mu$V        & 6,984 $\mu$V                  & \OK           & Filtre pondérateur désactivé                                         \\
    F.3         & Rapport signal sur bruit pondéré A, valeur calculée                    & 108 à 110 dB             & L 112,2 dB\newline R 112,4 dB & \OK           & Calculé au tableau~\ref{tab:ann_ap_bruit}                            \\
    F.4         & Plage dynamique AES17, 1 kHz à $-60$ dBFS, pondérée A, valeur calculée & 108 à 110 dB             & L 109,8 dB\newline R 110,1 dB & \OK           & Calculée au tableau~\ref{tab:ann_ap_bruit}                           \\
\end{longtable}
}

\section{Relevés chiffrés de la session Audio Precision}
\label{ann:mesures_ap}

Les niveaux sont saisis en volts efficaces. La colonne \gls{dbr} est normalisée au relevé à 1 kHz. Le THD+N est saisi en dB pour le balayage en fréquence et converti en pour cent pour faciliter la comparaison avec les valeurs de datasheet. Toutes les mesures sont non chargées sur les 100 k$\Omega$ de l'analyseur.

\begin{table}[H]
    \centering
    \caption{Réponse en fréquence et THD+N, source à 44,1 kHz et $-1$ dBFS, sortie à vide}
    \label{tab:ann_ap_reponse}
    \setlength{\tabcolsep}{4pt}
    \begin{tabular}{@{}r rr rr rr rr@{}}
        \toprule
                        & \multicolumn{2}{c}{\textbf{Niveau [V]}} & \multicolumn{2}{c}{\textbf{Écart [dBr]}} & \multicolumn{2}{c}{\textbf{THD+N [dB]}} & \multicolumn{2}{c}{\textbf{THD+N [\%]}}                                                     \\
        \cmidrule(lr){2-3} \cmidrule(lr){4-5} \cmidrule(lr){6-7} \cmidrule(lr){8-9}
        \textbf{f [Hz]} & \textbf{A}                              & \textbf{B}                               & \textbf{A}                              & \textbf{B}                              & \textbf{A} & \textbf{B} & \textbf{A} & \textbf{B} \\
        \midrule
        20              & 1,902                                   & 1,907                                    & $-0,018$                                & $-0,014$                                & $-93,6$    & $-93,5$    & 0,00209    & 0,00211    \\
        30              & 1,902                                   & 1,906                                    & $-0,018$                                & $-0,018$                                & $-93,5$    & $-93,7$    & 0,00211    & 0,00207    \\
        50              & 1,901                                   & 1,905                                    & $-0,023$                                & $-0,023$                                & $-94,2$    & $-93,9$    & 0,00195    & 0,00202    \\
        100             & 1,901                                   & 1,905                                    & $-0,023$                                & $-0,023$                                & $-93,7$    & $-93,9$    & 0,00207    & 0,00202    \\
        300             & 1,902                                   & 1,906                                    & $-0,018$                                & $-0,018$                                & $-93,7$    & $-94,0$    & 0,00207    & 0,00200    \\
        500             & 1,903                                   & 1,907                                    & $-0,014$                                & $-0,014$                                & $-94,0$    & $-94,3$    & 0,00200    & 0,00193    \\
        1000            & 1,906                                   & 1,910                                    & 0,000                                   & 0,000                                   & $-94,3$    & $-94,4$    & 0,00193    & 0,00191    \\
        3000            & 1,904                                   & 1,908                                    & $-0,009$                                & $-0,009$                                & $-94,0$    & $-94,1$    & 0,00200    & 0,00197    \\
        5000            & 1,905                                   & 1,909                                    & $-0,005$                                & $-0,005$                                & $-92,6$    & $-93,1$    & 0,00234    & 0,00221    \\
        10000           & 1,912                                   & 1,916                                    & 0,027                                   & 0,027                                   & $-93,0$    & $-93,9$    & 0,00224    & 0,00202    \\
        15000           & 1,914                                   & 1,918                                    & 0,036                                   & 0,036                                   & $-95,7$    & $-96,4$    & 0,00164    & 0,00151    \\
        18000           & 1,918                                   & 1,922                                    & 0,055                                   & 0,054                                   & $-94,5$    & $-95,2$    & 0,00188    & 0,00174    \\
        20000           & 1,915                                   & 1,919                                    & 0,041                                   & 0,041                                   & $-93,1$    & $-93,6$    & 0,00221    & 0,00209    \\
        \bottomrule
    \end{tabular}
\end{table}

% Courbes produites par script Python à partir du même jeu de données
\fig[H, width=1\textwidth]{Réponse en fréquence de la sortie jack, source à 44,1 kHz}{reponse_frequence.png}
\fig[H, width=1\textwidth]{THD+N en fonction de la fréquence, source à 44,1 kHz}{thdn_frequence.png}

L'écart maximal de niveau sur toute la bande audio est de 0,073 dB entre le point le plus bas, à 50 Hz et 100 Hz, et le point le plus haut, à 18 kHz. Le critère de $\pm$ 0,1 dB est donc tenu sans marge critique. La légère remontée au-delà de 10 kHz est cohérente avec la réponse du filtre numérique de reconstruction du \gls{DAC} et n'est pas imputable au filtre RC, dont la transparence est établie à la section~\ref{ann:mesures_filtre}.

\begin{table}[H]
    \centering
    \caption{THD+N en fonction du niveau, 1 kHz, sortie à vide}
    \label{tab:ann_ap_niveau}
    \small
    \setlength{\tabcolsep}{6pt}
    \begin{tabular}{@{}r rr rr@{}}
        \toprule
                               & \multicolumn{2}{c}{\textbf{Niveau [V]}} & \multicolumn{2}{c}{\textbf{THD+N [dB]}}                           \\
        \cmidrule(lr){2-3} \cmidrule(lr){4-5}
        \textbf{Niveau [dBFS]} & \textbf{A}                              & \textbf{B}                              & \textbf{A} & \textbf{B} \\
        \midrule
        0                      & 2,138                                   & 2,143                                   & $-94,8$    & $-95,2$    \\
        $-1$                   & 1,906                                   & 1,910                                   & $-94,6$    & $-94,9$    \\
        $-3$                   & 1,514                                   & 1,517                                   & $-93,9$    & $-94,9$    \\
        $-6$                   & 1,072                                   & 1,074                                   & $-92,5$    & $-93,4$    \\
        $-10$                  & 0,6761                                  & 0,6776                                  & $-90,0$    & $-91,0$    \\
        $-20$                  & 0,2136                                  & 0,2141                                  & $-81,6$    & $-83,0$    \\
        $-40$                  & 0,02136                                 & 0,02142                                 & $-61,7$    & $-63,6$    \\
        $-60$                  & 0,002137                                & 0,002142                                & $-41,7$    & $-43,6$    \\
        \bottomrule
    \end{tabular}
\end{table}

\fig[H, width=1\textwidth]{THD+N en fonction du niveau numérique, 1 kHz}{thdn_niveau.png}

En dessous de $-20$ dBFS, le THD+N décroît d'environ 20 dB par décade de niveau, signature d'un résidu dominé par le bruit et non par la distorsion. Au-dessus de $-6$ dBFS, la valeur se stabilise autour de $-95$ dB, ce qui correspond au plancher de distorsion du \gls{DAC}. Le croisement entre les deux régimes se situe entre $-10$ et $-6$ dBFS.

\begin{table}[H]
    \centering
    \caption{Bruit, rapport signal sur bruit et plage dynamique}
    \label{tab:ann_ap_bruit}
    \small
    \setlength{\tabcolsep}{6pt}
    \begin{tabular}{@{}>{\raggedright\arraybackslash}p{0.44\textwidth} r r l@{}}
        \toprule
        \textbf{Grandeur}                               & \textbf{Voie A} & \textbf{Voie B} & \textbf{Unité} \\
        \midrule
        Niveau de référence, 1 kHz à 0 dBFS             & 2,138           & 2,142           & V              \\
        Niveau mesuré à $-60$ dBFS                      & 0,002136        & 0,002141        & V              \\
        Bruit non pondéré, 22 Hz à 20 kHz               & 6,984           & 6,954           & $\mu$V         \\
        Bruit pondéré A                                 & 5,252           & 5,165           & $\mu$V         \\
        THD+N pondéré A pendant le signal à $-60$ dBFS  & $-49,8$         & $-50,1$         & dB             \\
        Résidu pondéré A pendant le signal à $-60$ dBFS & 6,912           & 6,693           & $\mu$V         \\
        Rapport signal sur bruit non pondéré            & 109,7           & 109,8           & dB             \\
        Rapport signal sur bruit pondéré A              & 112,2           & 112,4           & dB             \\
        Plage dynamique AES17, pondérée A               & 109,8           & 110,1           & dB             \\
        \bottomrule
    \end{tabular}
\end{table}
